\documentclass{beamer}

\usetheme{LMU}

\usepackage[utf8]{inputenc}
\usepackage[T1]{fontenc}

\usepackage[english]{babel}
\usepackage{url}
\usepackage{listings}
%\usepackage{cite}
%\usepackage{natbib}
\usepackage[backend=bibtex,style=authoryear,dashed=false]{biblatex}
\addbibresource{../bib/eigene.bib}
\addbibresource{../bib/itip-refs.bib}
\addbibresource{../bib/other-refs.bib}
\renewcommand{\bibfont}{\normalfont\scriptsize}
\setlength{\bibhang}{3ex}

\usepackage{amssymb, amsmath, amsfonts, enumerate}
%\usepackage{bbold}
\newcommand\hmmax{0}
\usepackage{bm}
%\usepackage{dsfont}
\usepackage{pxfonts}
\usepackage{xcolor}

\usepackage{tikz}
\usetikzlibrary{%
   arrows,%
   calc,%
   fit,%
   patterns,%
   plotmarks,%
   shapes.geometric,%
   shapes.misc,%
   shapes.symbols,%
   shapes.arrows,%
   shapes.callouts,%
   shapes.multipart,%
   shapes.gates.logic.US,%
   shapes.gates.logic.IEC,%
   er,%
   automata,%
   backgrounds,%
   chains,%
   topaths,%
   trees,%
   petri,%
   mindmap,%
   matrix,%
   calendar,%
   folding,%
   fadings,%
   through,%
   patterns,%
   positioning,%
   scopes,%
   decorations.fractals,%
   decorations.shapes,%
   decorations.text,%
   decorations.pathmorphing,%
   decorations.pathreplacing,%
   decorations.footprints,%
   decorations.markings,%
   shadows,%
   spy}
%\usepackage{bbold}
\usepackage{hyperref}

\setbeamertemplate{blocks}[rounded][shadow=true]
\definecolor{lmugreen2}{RGB}{0,120,94}
%\definecolor{lmugreen2}{RGB}{0,148,64}
\definecolor{lmugreen}{RGB}{0,140,84}
\definecolor{unidurham}{RGB}{126,49,123}

\parindent0pt
\setlength{\unitlength}{1ex}
\setlength{\fboxsep}{0ex}

\def\then{{\color{lmugreen}$\rule[0.35ex]{2ex}{0.5ex}\!\!\!\blacktriangleright$}}
%\def\then{{\color{lmugreen}$\blacktriangleright\!\blacktriangleright$}}
%\def\then{{\color{lmugreen}$\blacktriangleright$}}
%\def\then{{\color{lmugreen}$\Rrightarrow$}}
%\def\then{{\color{lmugreen}$\rhd$}}
%\def\then{{\color{lmugreen}$\gg\!\!\!\!\!\gg$}}
%\def\then{{\color{lmugreen}${\mathbf{\gg}}$}}
\def\play{{\color{lmugreen}$\blacktriangleright$}}

\def\rthen{{\color{lmugreen}$\rule[0.35ex]{0.5ex}{0.95ex}\rule[0.35ex]{1.3ex}{0.5ex}\!\!\!\blacktriangleright$}}

\def\thenthen{{\color{lmugreen}$\blacktriangleleft\!\!\!\rule[0.35ex]{2ex}{0.5ex}\!\!\!\blacktriangleright$}}

\def\gplus{{\color{lmugreen}\rule[0.45ex]{1.4ex}{0.4ex}\hspace{-0.9ex}\rule[0.0ex]{0.4ex}{1.3ex}\hspace{0.5ex}}}
\def\gminus{{\color{lmugreen}\rule[0.45ex]{1.4ex}{0.4ex}}}


\def\blau#1{{\color{lmugreen2}#1}}
\def\rot#1{{\color{red}#1}}
\def\gruen#1{{\color{blue}#1}}
%\def\gruen#1{{\color{gray}#1}}

%%%%%%%%%%%%%%%%%%%%%%%%%%%%%%%%%%%%%%%%%
%% Definitions & shortcuts for thesis  %%
%%%%%%%%%%%%%%%%%%%%%%%%%%%%%%%%%%%%%%%%%

\def\pdc{prior-data conflict}

\newcommand{\reals}{\mathbb{R}}

\newcommand{\dd}{\,\mathrm{d}}

\newcommand{\mbf}[1]{\mathbf{#1}}

\newcommand{\X}{\mbf{X}}
\newcommand{\x}{\mbf{x}}

\def\yzr{\rot{\yz}}
\def\ynr{\rot{\yn}}
\def\byzr{\rot{\byz}}
\def\bynr{\rot{\byn}}
\def\yzor{\rot{y\uz_1}}
\def\yzjr{\rot{y\uz_j}}
\def\yzkr{\rot{y\uz_k}}
\def\yzlr{\rot{\yzl}}
\def\yzur{\rot{\yzu}}
\def\ynjr{\rot{y\un_j}}
\def\ynlr{\rot{\ynl}}
\def\ynur{\rot{\ynu}}
\def\yzjlr#1{\rot{\ul{y}\uz_#1}}
\def\yzjur#1{\rot{\ol{y}\uz_#1}}


\def\nzg{\gruen{\nz}}
\def\nng{\gruen{\nn}}
\def\nzlg{\gruen{\nzl}}
\def\nzug{\gruen{\nzu}}
\def\nnlg{\gruen{\nnl}}
\def\nnug{\gruen{\nnu}}

\def\psib{\blau{\psi}}
\def\bpsib{\blau{{b}(\psi)}}


% ------------ shading start
\newsavebox{\tempbox}
\newcommand\leftrightshading[3]{%
  \begin{tikzfadingfrompicture}[name=inputtext]
    \node [text=white] {#1};
  \end{tikzfadingfrompicture}
  \begin{lrbox}{\tempbox}%
    \begin{tikzpicture}
      \node [text=white,inner sep=0pt,outer sep=0pt] (textnode) {#1};
      \shade[path fading=inputtext,fit fading=false,left color=#2,right color=#3]
      (textnode.south west) rectangle (textnode.north east);
    \end{tikzpicture}%
  \end{lrbox}
  % Now we use the fading in another picture:
  \usebox\tempbox{}%
}
% ------------ shading end


%\def\PZc{\mathrm I\!\Pi\uz}
\def\PZc{\leftrightshading{$\mathrm I\!\Pi\uz$}{blue}{red}}
%\def\PNc{\PN}
\def\PNc{\leftrightshading{$\mathrm I\!\Pi\un$}{blue}{red}}


\title[Generalised Bayesian Inference \& Prior-Data Conflict]
{Generalised Bayesian Inference\\ under Prior-Data Conflict}

\author%[]
       {Gero Walter}

\institute{Department of Statistics\\
           Ludwig-Maximilians-Universit\"at M\"unchen (LMU)\\
                      {}%\{gero.walter; thomas\}@stat.uni-muenchen.de
}

\date{October 25th, 2013}


\titlegraphic{
%\begin{center}
\includegraphics[scale=0.032]{lmu_logos/lmu_massiv.png}
\raisebox{1.375cm}{
\includegraphics[scale=0.4]{lmu_logos/lmu_statistic.pdf}
}
%\end{center}
}

%---- Notes ---------------------
\setbeameroption{hide notes}
%\setbeameroption{show only notes}

\setbeamertemplate{note page}{%
 {\Huge \theframenumber}\\
 \insertnote%
}

\begin{document}


\frame{
\titlepage
}

% ---------- brace with label -----------
%\draw[decorate,decoration={brace,mirror,raise=5pt},blue]
%\draw[thick,decorate,decoration=brace] (2.2,-1.5) -- (2.2,-2.7) node[midway,right,align=left] %
%{motivating example:\\ {\color{lmugreen} common-cause failure}};

\frame{\frametitle{Overview}
\begin{tikzpicture}%[spy using outlines={rectangle, magnification=3, size=1.5cm, connect spies}]
\path (-1.5,-5) rectangle (10.5,2);
\node at (4.5,-3.4) {\parbox{0.9\textwidth}{%
\begin{enumerate}
\item Bayesian inference \& prior-data conflict %and Strong prior-data agreement
\item Generalised Bayesian inference with sets of priors\\ (joint work with Thomas Augustin and Frank Coolen)
\item Common-cause failure modeling\\ (joint work with Matthias Troffaes and Dana Kelly)
\end{enumerate}
}};
\end{tikzpicture}
}

\addtocounter{framenumber}{-1} 
\frame{\frametitle{Overview}
\begin{tikzpicture}%[spy using outlines={rectangle, magnification=3, size=1.5cm, connect spies}]
\path (-1.5,-5) rectangle (10.5,2);
\node[draw] at (0,0) {\includegraphics[scale=0.12]{../diss-toc-1.pdf}};
\node[draw] at (3,0) {\includegraphics[scale=0.12]{../diss-toc-2.pdf}};
\node[draw] at (6,0) {\includegraphics[scale=0.12]{../diss-toc-3.pdf}};
\node[draw] at (9,0) {\includegraphics[scale=0.12]{../diss-toc-4.pdf}};
%
\node at (4.5,-3.4) {\parbox{0.9\textwidth}{%
\begin{enumerate}
\item Bayesian inference \& prior-data conflict %and Strong prior-data agreement
\item Generalised Bayesian inference with sets of priors\\ (joint work with Thomas Augustin and Frank Coolen)
\item Common-cause failure modeling\\ (joint work with Matthias Troffaes and Dana Kelly)
\end{enumerate}
}};
\end{tikzpicture}
}



% page 1
\addtocounter{framenumber}{-1} 
\frame{\frametitle{Overview}
\begin{tikzpicture}%[spy using outlines={rectangle, magnification=3, size=1.5cm, connect spies}]
\path (-1.5,-5) rectangle (10.5,2);
\node[draw, yellow!60!red, fill=white!60!yellow] at (0,0) {\includegraphics[scale=0.12]{../diss-toc-1.pdf}};
\node[draw] at (3,0) {\includegraphics[scale=0.12]{../diss-toc-2.pdf}};
\node[draw] at (6,0) {\includegraphics[scale=0.12]{../diss-toc-3.pdf}};
\node[draw] at (9,0) {\includegraphics[scale=0.12]{../diss-toc-4.pdf}};
%
\node[draw, yellow!60!red, fill=white] at (1.3,-1.5) {\includegraphics[scale=0.225]{../diss-toc-1.pdf}};
\end{tikzpicture}
}

\addtocounter{framenumber}{-1} 
\frame{\frametitle{Overview}
\begin{tikzpicture}[spy using outlines={rectangle, magnification=3, width=11cm, connect spies}]
\path (-1.5,-5) rectangle (10.5,2);
\node[draw, yellow!60!red, fill=white!60!yellow] at (0,0) {\includegraphics[scale=0.12]{../diss-toc-1.pdf}};
\node[draw] at (3,0) {\includegraphics[scale=0.12]{../diss-toc-2.pdf}};
\node[draw] at (6,0) {\includegraphics[scale=0.12]{../diss-toc-3.pdf}};
\node[draw] at (9,0) {\includegraphics[scale=0.12]{../diss-toc-4.pdf}};
\node[draw, yellow!60!red, fill=white] at (1.3,-1.5) {\includegraphics[scale=0.225]{../diss-toc-1.pdf}};
\spy[red, height=0.55cm] on (1.34,-1.3) in node [fill=white] at (5, 1.5);
\end{tikzpicture}
}

\addtocounter{framenumber}{-1} 
\frame{\frametitle{Overview}
\begin{tikzpicture}[spy using outlines={rectangle, magnification=3, width=11cm, connect spies}]
\path (-1.5,-5) rectangle (10.5,2);
\node[draw, yellow!60!red, fill=white!60!yellow] at (0,0) {\includegraphics[scale=0.12]{../diss-toc-1.pdf}};
\node[draw] at (3,0) {\includegraphics[scale=0.12]{../diss-toc-2.pdf}};
\node[draw] at (6,0) {\includegraphics[scale=0.12]{../diss-toc-3.pdf}};
\node[draw] at (9,0) {\includegraphics[scale=0.12]{../diss-toc-4.pdf}};
\node[draw, yellow!60!red, fill=white] at (1.3,-1.5) {\includegraphics[scale=0.225]{../diss-toc-1.pdf}};
\spy[red, height=0.55cm] on (1.34,-1.3) in node [fill=white] at (5, 1.5);
\spy[red, height=3.55cm] on (1.34,-3.4) in node [fill=white] at (5,-0.8);
\end{tikzpicture}
}

% page 2
\addtocounter{framenumber}{-1} 
\frame{\frametitle{Overview}
\begin{tikzpicture}%[spy using outlines={rectangle, magnification=3, size=1.5cm, connect spies}]
\path (-1.5,-5) rectangle (10.5,2);
\node[draw] at (0,0) {\includegraphics[scale=0.12]{../diss-toc-1.pdf}};
\node[draw, yellow!60!red, fill=white!60!yellow] at (3,0) {\includegraphics[scale=0.12]{../diss-toc-2.pdf}};
\node[draw] at (6,0) {\includegraphics[scale=0.12]{../diss-toc-3.pdf}};
\node[draw] at (9,0) {\includegraphics[scale=0.12]{../diss-toc-4.pdf}};
%
\node[draw, yellow!60!red, fill=white] at (1.3,-1.5) {\includegraphics[scale=0.225]{../diss-toc-2.pdf}};
\end{tikzpicture}
}

\addtocounter{framenumber}{-1} 
\frame{\frametitle{Overview}
\begin{tikzpicture}[spy using outlines={rectangle, magnification=3, width=11cm}]
\path (-1.5,-5) rectangle (10.5,2);
\node[draw] at (0,0) {\includegraphics[scale=0.12]{../diss-toc-1.pdf}};
\node[draw, yellow!60!red, fill=white!60!yellow] at (3,0) {\includegraphics[scale=0.12]{../diss-toc-2.pdf}};
\node[draw] at (6,0) {\includegraphics[scale=0.12]{../diss-toc-3.pdf}};
\node[draw] at (9,0) {\includegraphics[scale=0.12]{../diss-toc-4.pdf}};
%
\node[draw, yellow!60!red, fill=white] at (1.3,-1.5) {\includegraphics[scale=0.225]{../diss-toc-2.pdf}};
\spy[red, height=0.55cm] on (1.34,-1.53) in node [fill=white] at (5,-0.4);
\end{tikzpicture}
}

\addtocounter{framenumber}{-1} 
\frame{\frametitle{Overview}
\begin{tikzpicture}[spy using outlines={rectangle, magnification=3, width=11cm}]
\path (-1.5,-5) rectangle (10.5,2);
\node[draw] at (0,0) {\includegraphics[scale=0.12]{../diss-toc-1.pdf}};
\node[draw, yellow!60!red, fill=white!60!yellow] at (3,0) {\includegraphics[scale=0.12]{../diss-toc-2.pdf}};
\node[draw] at (6,0) {\includegraphics[scale=0.12]{../diss-toc-3.pdf}};
\node[draw] at (9,0) {\includegraphics[scale=0.12]{../diss-toc-4.pdf}};
%
\node[draw, yellow!60!red, fill=white] at (1.3,-1.5) {\includegraphics[scale=0.225]{../diss-toc-2.pdf}};
\spy[red, height=0.55cm] on (1.34,-1.53) in node [fill=white] at (5,-0.4);
\spy[red, height=0.55cm] on (1.34,-0.85) in node [fill=white] at (5, 0.3);
\end{tikzpicture}
}

\addtocounter{framenumber}{-1} 
\frame{\frametitle{Overview}
\begin{tikzpicture}[spy using outlines={rectangle, magnification=3, width=11cm}]
\path (-1.5,-5) rectangle (10.5,2);
\node[draw] at (0,0) {\includegraphics[scale=0.12]{../diss-toc-1.pdf}};
\node[draw, yellow!60!red, fill=white!60!yellow] at (3,0) {\includegraphics[scale=0.12]{../diss-toc-2.pdf}};
\node[draw] at (6,0) {\includegraphics[scale=0.12]{../diss-toc-3.pdf}};
\node[draw] at (9,0) {\includegraphics[scale=0.12]{../diss-toc-4.pdf}};
%
\node[draw, yellow!60!red, fill=white] at (1.3,-1.5) {\includegraphics[scale=0.225]{../diss-toc-2.pdf}};
\spy[red, height=0.55cm] on (1.34,-1.53) in node [fill=white] at (5,-0.4);
\spy[red, height=0.55cm] on (1.34,-0.85) in node [fill=white] at (5, 0.3);
\spy[red, height=1.05cm] on (1.34, 0.78) in node [fill=white] at (5, 1.3);
\end{tikzpicture}
}

\addtocounter{framenumber}{-1} 
\frame{\frametitle{Overview}
\begin{tikzpicture}[spy using outlines={rectangle, magnification=3, width=11cm}]
\path (-1.5,-5) rectangle (10.5,2);
\node[draw] at (0,0) {\includegraphics[scale=0.12]{../diss-toc-1.pdf}};
\node[draw, yellow!60!red, fill=white!60!yellow] at (3,0) {\includegraphics[scale=0.12]{../diss-toc-2.pdf}};
\node[draw] at (6,0) {\includegraphics[scale=0.12]{../diss-toc-3.pdf}};
\node[draw] at (9,0) {\includegraphics[scale=0.12]{../diss-toc-4.pdf}};
%
\node[draw, yellow!60!red, fill=white] at (1.3,-1.5) {\includegraphics[scale=0.225]{../diss-toc-2.pdf}};
\spy[red, height=0.55cm] on (1.34,-1.53) in node [fill=white] at (5,-0.4);
\spy[red, height=0.55cm] on (1.34,-0.85) in node [fill=white] at (5, 0.3);
\spy[red, height=1.05cm] on (1.34, 0.78) in node [fill=white] at (5, 1.3);
\spy[red, height=2.45cm] on (1.34,-2.43) in node [fill=white] at (5,-3.7);
\end{tikzpicture}
}

% page 3
\addtocounter{framenumber}{-1} 
\frame{\frametitle{Overview}
\begin{tikzpicture}%[spy using outlines={rectangle, magnification=3, size=1.5cm, connect spies}]
\path (-1.5,-5) rectangle (10.5,2);
\node[draw] at (0,0) {\includegraphics[scale=0.12]{../diss-toc-1.pdf}};
\node[draw] at (3,0) {\includegraphics[scale=0.12]{../diss-toc-2.pdf}};
\node[draw, yellow!60!red, fill=white!60!yellow] at (6,0) {\includegraphics[scale=0.12]{../diss-toc-3.pdf}};
\node[draw] at (9,0) {\includegraphics[scale=0.12]{../diss-toc-4.pdf}};
%
\node[draw, yellow!60!red, fill=white] at (1.3,-1.5) {\includegraphics[scale=0.225]{../diss-toc-3.pdf}};
\end{tikzpicture}
}

\addtocounter{framenumber}{-1} 
\frame{\frametitle{Overview}
\begin{tikzpicture}[spy using outlines={rectangle, magnification=3, width=11cm}]
\path (-1.5,-5) rectangle (10.5,2);
\node[draw] at (0,0) {\includegraphics[scale=0.12]{../diss-toc-1.pdf}};
\node[draw] at (3,0) {\includegraphics[scale=0.12]{../diss-toc-2.pdf}};
\node[draw, yellow!60!red, fill=white!60!yellow] at (6,0) {\includegraphics[scale=0.12]{../diss-toc-3.pdf}};
\node[draw] at (9,0) {\includegraphics[scale=0.12]{../diss-toc-4.pdf}};
%
\node[draw, yellow!60!red, fill=white] at (1.3,-1.5) {\includegraphics[scale=0.225]{../diss-toc-3.pdf}};
\spy[red, height=2.90cm] on (1.34, 0.36) in node [fill=white] at (5,-2.1);
\end{tikzpicture}
}

\addtocounter{framenumber}{-1} 
\frame{\frametitle{Overview}
\begin{tikzpicture}[spy using outlines={rectangle, magnification=3, width=11cm}]
\path (-1.5,-5) rectangle (10.5,2);
\node[draw] at (0,0) {\includegraphics[scale=0.12]{../diss-toc-1.pdf}};
\node[draw] at (3,0) {\includegraphics[scale=0.12]{../diss-toc-2.pdf}};
\node[draw, yellow!60!red, fill=white!60!yellow] at (6,0) {\includegraphics[scale=0.12]{../diss-toc-3.pdf}};
\node[draw] at (9,0) {\includegraphics[scale=0.12]{../diss-toc-4.pdf}};
%
\node[draw, yellow!60!red, fill=white] at (1.3,-1.5) {\includegraphics[scale=0.225]{../diss-toc-3.pdf}};
\spy[red, height=2.90cm] on (1.34, 0.36) in node [fill=white] at (5,-2.1);
\spy[red, height=0.55cm] on (1.34,-0.28) in node [fill=white] at (5,-4.0);
\end{tikzpicture}
}

% page 4
\addtocounter{framenumber}{-1} 
\frame{\frametitle{Overview}
\begin{tikzpicture}%[spy using outlines={rectangle, magnification=3, size=1.5cm, connect spies}]
\path (-1.5,-5) rectangle (10.5,2);
\node[draw] at (0,0) {\includegraphics[scale=0.12]{../diss-toc-1.pdf}};
\node[draw] at (3,0) {\includegraphics[scale=0.12]{../diss-toc-2.pdf}};
\node[draw] at (6,0) {\includegraphics[scale=0.12]{../diss-toc-3.pdf}};
\node[draw, yellow!60!red, fill=white!60!yellow] at (9,0) {\includegraphics[scale=0.12]{../diss-toc-4.pdf}};
%
\node[draw, yellow!60!red, fill=white] at (1.3,-1.5) {\includegraphics[scale=0.225]{../diss-toc-4.pdf}};
\end{tikzpicture}
}

\addtocounter{framenumber}{-1} 
\frame{\frametitle{Overview}
\begin{tikzpicture}[spy using outlines={rectangle, magnification=3, width=11cm}]
\path (-1.5,-5) rectangle (10.5,2);
\node[draw] at (0,0) {\includegraphics[scale=0.12]{../diss-toc-1.pdf}};
\node[draw] at (3,0) {\includegraphics[scale=0.12]{../diss-toc-2.pdf}};
\node[draw] at (6,0) {\includegraphics[scale=0.12]{../diss-toc-3.pdf}};
\node[draw, yellow!60!red, fill=white!60!yellow] at (9,0) {\includegraphics[scale=0.12]{../diss-toc-4.pdf}};
%
\node[draw, yellow!60!red, fill=white] at (1.3,-1.5) {\includegraphics[scale=0.225]{../diss-toc-4.pdf}};
\spy[red, height=0.85cm] on (1.34, 0.47) in node [fill=white] at (5,-3.0);
\end{tikzpicture}
}

\iffalse
% ----------------------------------------fast-forward compiling
\section{Common-cause Failure Modelling}
\subsection{Alpha-Factor Model}
\subsection{Bayesian Inference}
\subsection{Imprecise Dirichlet Model}
\frame{
Outline of the talk***
\begin{enumerate}
\item Bayesian inference \& prior-data conflict %and Strong prior-data agreement
\item Generalised Bayesian inference with sets of priors\\ (joint work with Thomas Augustin and Frank Coolen)
\item Common-cause failure modeling\\ (joint work with Matthias Troffaes and Dana Kelly)
\end{enumerate}
}
\end{document}
% ----------------------------------------fast-forward compiling
\fi


\iffalse
\frame{\frametitle{Overview}
\begin{enumerate}
\item Introduction
 \begin{enumerate}
 \item Preliminaries
 \item Some Fundamentals
 \item Dirichlet-Multinomial Model for Common-Cause Failure
 \end{enumerate}
\item Imprecise Probability as Foundation of Generalised Bayesian Inference
 \begin{enumerate}
 \item Imprecise or Interval Probability
 \item Motives for the Use of Imprecise Probability
 \end{enumerate}
\item Generalised Bayesian Inference with Sets of Conjugate Priors in Exponential Families
 \begin{enumerate}
 \item Model Overview and Discussion
 \item Alternative Models Using Sets of Priors
 \item Imprecision and Prior-Data Conflict in Generalised Bayesian Inference
 \item The \texttt{luck} Package
 \item On Prior-Data Conflict inPredictive Bernoulli Inferences
 \end{enumerate}
\end{enumerate}
}

\frame{\frametitle{Overview}
\begin{enumerate}
\item[4.] Concluding Remarks
 \begin{enumerate}
 \item[4.1] Summary
 \item[4.2] Discussion
 \item[4.3] Outlook
 \end{enumerate}
\item[A.] Appendix
 \begin{enumerate}
 \item[A.1] Bayesian Linear Regression: Different Conjugate Models and Their (In)Sensitivity to Prior-Data Conflict
 \item[A.2] A Parameter Set Shape for Strong Prior-Data Agreement Modelling
 \end{enumerate}
\end{enumerate}

Outline of the talk***
\begin{enumerate}
\item Bayesian inference \& prior-data conflict %and Strong prior-data agreement
\item Generalised Bayesian inference with sets of priors\\ (joint work with Thomas Augustin and Frank Coolen)
\item Common-cause failure modeling\\ (joint work with Matthias Troffaes and Dana Kelly)
\end{enumerate}
}
\fi

\section{Bayesian Inference \& Prior-Data Conflict}

\subsection{Bayesian Inference}

\frame{\frametitle{Bayesian Inference \& Prior-Data Conflict}

\uncover<1->{%
\begin{block}{The Bayesian approach to statistical inference}
\centerline{prior $p(\vartheta)$ + likelihood $f(\x\mid\vartheta)$\hspace*{1ex} \then\ \ posterior $p(\vartheta\mid\x)$}
%\quad using Bayes' Rule}
All inferences are based on the posterior (e.g., point estimate, \ldots)
\end{block}}
\uncover<2->{%
Assigning a certain prior distribution on $\vartheta$\\
\hfill $=$ defining a conglomerate of probability statements (on $\vartheta$)
}
\uncover<3->{%
\begin{block}{Prior-Data Conflict}
\begin{itemize}
\item \emph{informative prior beliefs} and \emph{trusted data}\\ %\rule{0ex}{3ex}\\
(sampling model correct, no outliers, etc.) are in conflict%\\%[2ex]
\item ``[\ldots] the prior [places] its mass primarily on distributions
in the sampling model for which the observed data is surprising''\\
(Evans \& Moshonov, 2006)
\item there are not enough data to overrule the prior
\end{itemize}
\end{block}}
%\note<3>[item]{real-world example where \pdc\ is relevant\\ at the end of the talk}
}


\frame{\frametitle{Prior-Data Conflict: Basic Example}

%Example Beta-Binom, show point, line, rectangle
\begin{itemize}[<+->]
\item Bernoulli observations: 0/1 observations (team wins no/yes)
\item given: a set of observations and strong prior information
%\item given: a set of observations (team won 12 out of 16 matches)
%\item additional to observations, we have strong prior information\\ (we are convinced that $P(\text{win})$ should be around $0.75$)
\item we are, e.g., interested in (predictive) probability $P$ that team wins in the next match
%(or some other inferences) %(predictive probability!)
%\item standard statistical model for this situation: Beta-Bernoulli/Binomial Model
%\item \textbf{\pdc:} if $P(\text{heads})$ for the coin is actually very different from our prior guess
%(i.e., prior information and data are in conflict), this should show up in the %predictive
%inferences (probability $\mbox{P}$ and, e.g., confidence intervals)
\end{itemize}
\uncover<4->{
\begin{block}{Beta-Binomial Model}
\begin{tabular}{r|lcl}
data :           & $s\mid \theta$          & $\sim$ & $\bin(n,\theta)$   \\[0.5ex]
conjugate prior: & $\theta\mid \nzg, \yzr$ & $\sim$ & $\be(\nzg,\, \yzr)$  \\[0.5ex]
\cline{1-4}
posterior:       & $\theta\mid \nng, \ynr$ & $\sim$ & $\be(\nng,\, \ynr)$\rule{0ex}{2.5ex}
%\quad ($\frac{\tau(\x)}{n} = \frac{s}{n}$)\rule{0ex}{2.5ex}
\end{tabular}
\end{block}
where $s$ = number of wins in the $n$ matches observed
}
\note<4>[item]{why conjugate priors? easy to calculate, posterior tractable (known family)}
}

\subsection{Example: Beta-Binomial Model}

\frame{\frametitle{Beta-Binomial Model (BBM)}

\uncover<1->{
\begin{block}{Beta-Binomial Model}
\begin{tabular}{r|lcl}
data :           & $s\mid \theta$          & $\sim$ & $\bin(n,\theta)$   \\[0.5ex]
conjugate prior: & $\theta\mid \nzg, \yzr$ & $\sim$ & $\be(\nzg,\, \yzr)$  \\[0.5ex]
\cline{1-4}
posterior:       & $\theta\mid \nng, \ynr$ & $\sim$ & $\be(\nng,\, \ynr)$\rule{0ex}{2.5ex}
%\quad ($\frac{\tau(\x)}{n} = \frac{s}{n}$)\rule{0ex}{2.5ex}
\end{tabular}
\end{block}
}
\begin{align*}
\uncover<2->{%
P = \E[\theta\mid\nng,\ynr]
}%
\uncover<3->{%
= \rot{\yn}
= \frac{\gruen{\nz}}{\gruen{\nz} + n} \cdot \rot{\yz} + \frac{n}{\gruen{\nz} + n} \cdot \frac{s}{n}
}%
\end{align*}\vspace*{-3ex}
\begin{align*}
\uncover<4->{%
\gruen{\nn} &= \gruen{\nz} + n
}%
\uncover<5->{%
& \V(\theta\mid\nng,\ynr) &= \frac{\rot{\yn}(1-\rot{\yn})}{\gruen{\nn} + 1}
}
\end{align*}

}

\frame{\frametitle{Beta-Binomial Model (BBM)}

\hspace*{-12ex}
\begin{columns}%[T]
\begin{column}{0.55\textwidth}
\begin{tikzpicture}
\pgftransformscale{0.025}
\uncover<1>{
% Created by tikzDevice version 0.5.0 on 2011-07-18 16:47:19
\begin{scope}
\path[clip] ( 55.20, 49.20) rectangle (250.54,221.74);
\definecolor[named]{drawColor}{rgb}{0.82,0.27,0.19}
\definecolor[named]{fillColor}{rgb}{0.88,0.08,0.52}
\definecolor[named]{drawColor}{rgb}{0.00,0.00,0.00}
\definecolor[named]{fillColor}{rgb}{0.00,0.00,0.00}

\draw[color=drawColor,line cap=round,line join=round,fill=fillColor,] ( 96.89,175.41) circle (  2.25);
\end{scope}
\begin{scope}
\path[clip] (  0.00,  0.00) rectangle (252.94,252.94);
\definecolor[named]{drawColor}{rgb}{0.82,0.27,0.19}
\definecolor[named]{fillColor}{rgb}{0.88,0.08,0.52}
\definecolor[named]{drawColor}{rgb}{0.00,0.00,0.00}

\draw[color=drawColor,line cap=round,line join=round,fill opacity=0.00,] ( 71.05, 49.20) -- (243.31, 49.20);

\draw[color=drawColor,line cap=round,line join=round,fill opacity=0.00,] ( 71.05, 49.20) -- ( 71.05, 43.20);

\draw[color=drawColor,line cap=round,line join=round,fill opacity=0.00,] (114.11, 49.20) -- (114.11, 43.20);

\draw[color=drawColor,line cap=round,line join=round,fill opacity=0.00,] (157.18, 49.20) -- (157.18, 43.20);

\draw[color=drawColor,line cap=round,line join=round,fill opacity=0.00,] (200.24, 49.20) -- (200.24, 43.20);

\draw[color=drawColor,line cap=round,line join=round,fill opacity=0.00,] (243.31, 49.20) -- (243.31, 43.20);

\node[color=drawColor,anchor=base,inner sep=0pt, outer sep=0pt, scale=  0.90] at ( 71.05, 25.20) {5%
};

\node[color=drawColor,anchor=base,inner sep=0pt, outer sep=0pt, scale=  0.90] at (114.11, 25.20) {10%
};

\node[color=drawColor,anchor=base,inner sep=0pt, outer sep=0pt, scale=  0.90] at (157.18, 25.20) {15%
};

\node[color=drawColor,anchor=base,inner sep=0pt, outer sep=0pt, scale=  0.90] at (200.24, 25.20) {20%
};

\node[color=drawColor,anchor=base,inner sep=0pt, outer sep=0pt, scale=  0.90] at (243.31, 25.20) {25%
};

\draw[color=drawColor,line cap=round,line join=round,fill opacity=0.00,] ( 55.20, 55.59) -- ( 55.20,215.35);

\draw[color=drawColor,line cap=round,line join=round,fill opacity=0.00,] ( 55.20, 55.59) -- ( 49.20, 55.59);

\draw[color=drawColor,line cap=round,line join=round,fill opacity=0.00,] ( 55.20, 87.54) -- ( 49.20, 87.54);

\draw[color=drawColor,line cap=round,line join=round,fill opacity=0.00,] ( 55.20,119.50) -- ( 49.20,119.50);

\draw[color=drawColor,line cap=round,line join=round,fill opacity=0.00,] ( 55.20,151.45) -- ( 49.20,151.45);

\draw[color=drawColor,line cap=round,line join=round,fill opacity=0.00,] ( 55.20,183.40) -- ( 49.20,183.40);

\draw[color=drawColor,line cap=round,line join=round,fill opacity=0.00,] ( 55.20,215.35) -- ( 49.20,215.35);

\node[rotate= 90.00,color=drawColor,anchor=base,inner sep=0pt, outer sep=0pt, scale=  0.90] at ( 43.20, 55.59) {0.0%
};

\node[rotate= 90.00,color=drawColor,anchor=base,inner sep=0pt, outer sep=0pt, scale=  0.90] at ( 43.20, 87.54) {0.2%
};

\node[rotate= 90.00,color=drawColor,anchor=base,inner sep=0pt, outer sep=0pt, scale=  0.90] at ( 43.20,119.50) {0.4%
};

\node[rotate= 90.00,color=drawColor,anchor=base,inner sep=0pt, outer sep=0pt, scale=  0.90] at ( 43.20,151.45) {0.6%
};

\node[rotate= 90.00,color=drawColor,anchor=base,inner sep=0pt, outer sep=0pt, scale=  0.90] at ( 43.20,183.40) {0.8%
};

\node[rotate= 90.00,color=drawColor,anchor=base,inner sep=0pt, outer sep=0pt, scale=  0.90] at ( 43.20,215.35) {1.0%
};

\draw[color=drawColor,line cap=round,line join=round,fill opacity=0.00,] ( 55.20, 49.20) --
    (250.54, 49.20) --
    (250.54,221.74) --
    ( 55.20,221.74) --
    ( 55.20, 49.20);
\end{scope}
\begin{scope}
\path[clip] (  0.00,  0.00) rectangle (252.94,252.94);
\definecolor[named]{drawColor}{rgb}{0.82,0.27,0.19}
\definecolor[named]{fillColor}{rgb}{0.88,0.08,0.52}
\definecolor[named]{drawColor}{rgb}{0.00,0.00,0.00}

\node[color=drawColor,anchor=base,inner sep=0pt, outer sep=0pt, scale=  1.00] at (152.87,  1.20) {$\gruen{\nz}$ resp. $\gruen{\nn}$%
};

\node[rotate= 90.00,color=drawColor,anchor=base,inner sep=0pt, outer sep=0pt, scale=  1.00] at ( 19.20,135.47) {$\rot{\yz}$ resp. $\rot{\yn}$%
};
\end{scope}

}
\uncover<2>{
% Created by tikzDevice version 0.5.0 on 2011-07-18 12:18:44
\begin{scope}
\path[clip] ( 55.20, 49.20) rectangle (250.54,221.74);
\definecolor[named]{fillColor}{rgb}{0.88,0.08,0.52}
\definecolor[named]{drawColor}{rgb}{0.00,0.00,0.00}
\definecolor[named]{fillColor}{rgb}{0.00,0.00,0.00}

\draw[color=drawColor,line cap=round,line join=round,fill=fillColor,] ( 96.89,175.41) circle (  2.25);

\draw[color=drawColor,line cap=round,line join=round,fill=fillColor,] (234.70,175.41) circle (  2.25);
\end{scope}
\begin{scope}
\path[clip] (  0.00,  0.00) rectangle (252.94,252.94);
\definecolor[named]{fillColor}{rgb}{0.88,0.08,0.52}
\definecolor[named]{drawColor}{rgb}{0.00,0.00,0.00}

\draw[color=drawColor,line cap=round,line join=round,fill opacity=0.00,] ( 71.05, 49.20) -- (243.31, 49.20);

\draw[color=drawColor,line cap=round,line join=round,fill opacity=0.00,] ( 71.05, 49.20) -- ( 71.05, 43.20);

\draw[color=drawColor,line cap=round,line join=round,fill opacity=0.00,] (114.11, 49.20) -- (114.11, 43.20);

\draw[color=drawColor,line cap=round,line join=round,fill opacity=0.00,] (157.18, 49.20) -- (157.18, 43.20);

\draw[color=drawColor,line cap=round,line join=round,fill opacity=0.00,] (200.24, 49.20) -- (200.24, 43.20);

\draw[color=drawColor,line cap=round,line join=round,fill opacity=0.00,] (243.31, 49.20) -- (243.31, 43.20);

\node[color=drawColor,anchor=base,inner sep=0pt, outer sep=0pt, scale=  0.90] at ( 71.05, 25.20) {5%
};

\node[color=drawColor,anchor=base,inner sep=0pt, outer sep=0pt, scale=  0.90] at (114.11, 25.20) {10%
};

\node[color=drawColor,anchor=base,inner sep=0pt, outer sep=0pt, scale=  0.90] at (157.18, 25.20) {15%
};

\node[color=drawColor,anchor=base,inner sep=0pt, outer sep=0pt, scale=  0.90] at (200.24, 25.20) {20%
};

\node[color=drawColor,anchor=base,inner sep=0pt, outer sep=0pt, scale=  0.90] at (243.31, 25.20) {25%
};

\draw[color=drawColor,line cap=round,line join=round,fill opacity=0.00,] ( 55.20, 55.59) -- ( 55.20,215.35);

\draw[color=drawColor,line cap=round,line join=round,fill opacity=0.00,] ( 55.20, 55.59) -- ( 49.20, 55.59);

\draw[color=drawColor,line cap=round,line join=round,fill opacity=0.00,] ( 55.20, 87.54) -- ( 49.20, 87.54);

\draw[color=drawColor,line cap=round,line join=round,fill opacity=0.00,] ( 55.20,119.50) -- ( 49.20,119.50);

\draw[color=drawColor,line cap=round,line join=round,fill opacity=0.00,] ( 55.20,151.45) -- ( 49.20,151.45);

\draw[color=drawColor,line cap=round,line join=round,fill opacity=0.00,] ( 55.20,183.40) -- ( 49.20,183.40);

\draw[color=drawColor,line cap=round,line join=round,fill opacity=0.00,] ( 55.20,215.35) -- ( 49.20,215.35);

\node[rotate= 90.00,color=drawColor,anchor=base,inner sep=0pt, outer sep=0pt, scale=  0.90] at ( 43.20, 55.59) {0.0%
};

\node[rotate= 90.00,color=drawColor,anchor=base,inner sep=0pt, outer sep=0pt, scale=  0.90] at ( 43.20, 87.54) {0.2%
};

\node[rotate= 90.00,color=drawColor,anchor=base,inner sep=0pt, outer sep=0pt, scale=  0.90] at ( 43.20,119.50) {0.4%
};

\node[rotate= 90.00,color=drawColor,anchor=base,inner sep=0pt, outer sep=0pt, scale=  0.90] at ( 43.20,151.45) {0.6%
};

\node[rotate= 90.00,color=drawColor,anchor=base,inner sep=0pt, outer sep=0pt, scale=  0.90] at ( 43.20,183.40) {0.8%
};

\node[rotate= 90.00,color=drawColor,anchor=base,inner sep=0pt, outer sep=0pt, scale=  0.90] at ( 43.20,215.35) {1.0%
};

\draw[color=drawColor,line cap=round,line join=round,fill opacity=0.00,] ( 55.20, 49.20) --
    (250.54, 49.20) --
    (250.54,221.74) --
    ( 55.20,221.74) --
    ( 55.20, 49.20);
\end{scope}
\begin{scope}
\path[clip] (  0.00,  0.00) rectangle (252.94,252.94);
\definecolor[named]{fillColor}{rgb}{0.88,0.08,0.52}
\definecolor[named]{drawColor}{rgb}{0.00,0.00,0.00}

\node[color=drawColor,anchor=base,inner sep=0pt, outer sep=0pt, scale=  1.00] at (152.87,  1.20) {$\gruen{\nz}$ resp. $\gruen{\nn}$%
};

\node[rotate= 90.00,color=drawColor,anchor=base,inner sep=0pt, outer sep=0pt, scale=  1.00] at ( 19.20,135.47) {$\rot{\yz}$ resp. $\rot{\yn}$%
};
\end{scope}

\draw[-stealth,very thick] (110,175) -- (220,175) node [above,midway] {12 out of 16};
}
\uncover<3->{
% Created by tikzDevice version 0.5.0 on 2011-07-18 12:18:48
\begin{scope}
\path[clip] ( 55.20, 49.20) rectangle (250.54,221.74);
\definecolor[named]{drawColor}{rgb}{0.60,0.27,0.19}
\definecolor[named]{fillColor}{rgb}{0.88,0.08,0.52}
\definecolor[named]{drawColor}{rgb}{0.00,0.00,0.00}
\definecolor[named]{fillColor}{rgb}{0.00,0.00,0.00}

\draw[color=drawColor,line cap=round,line join=round,fill=fillColor,] ( 96.89, 95.53) circle (  2.25);

\draw[color=drawColor,line cap=round,line join=round,fill=fillColor,] ( 96.89,175.41) circle (  2.25);

\draw[color=drawColor,line cap=round,line join=round,fill=fillColor,] (234.70,175.41) circle (  2.25);
\end{scope}
\begin{scope}
\path[clip] (  0.00,  0.00) rectangle (252.94,252.94);
\definecolor[named]{drawColor}{rgb}{0.60,0.27,0.19}
\definecolor[named]{fillColor}{rgb}{0.88,0.08,0.52}
\definecolor[named]{drawColor}{rgb}{0.00,0.00,0.00}

\draw[color=drawColor,line cap=round,line join=round,fill opacity=0.00,] ( 71.05, 49.20) -- (243.31, 49.20);

\draw[color=drawColor,line cap=round,line join=round,fill opacity=0.00,] ( 71.05, 49.20) -- ( 71.05, 43.20);

\draw[color=drawColor,line cap=round,line join=round,fill opacity=0.00,] (114.11, 49.20) -- (114.11, 43.20);

\draw[color=drawColor,line cap=round,line join=round,fill opacity=0.00,] (157.18, 49.20) -- (157.18, 43.20);

\draw[color=drawColor,line cap=round,line join=round,fill opacity=0.00,] (200.24, 49.20) -- (200.24, 43.20);

\draw[color=drawColor,line cap=round,line join=round,fill opacity=0.00,] (243.31, 49.20) -- (243.31, 43.20);

\node[color=drawColor,anchor=base,inner sep=0pt, outer sep=0pt, scale=  0.90] at ( 71.05, 25.20) {5%
};

\node[color=drawColor,anchor=base,inner sep=0pt, outer sep=0pt, scale=  0.90] at (114.11, 25.20) {10%
};

\node[color=drawColor,anchor=base,inner sep=0pt, outer sep=0pt, scale=  0.90] at (157.18, 25.20) {15%
};

\node[color=drawColor,anchor=base,inner sep=0pt, outer sep=0pt, scale=  0.90] at (200.24, 25.20) {20%
};

\node[color=drawColor,anchor=base,inner sep=0pt, outer sep=0pt, scale=  0.90] at (243.31, 25.20) {25%
};

\draw[color=drawColor,line cap=round,line join=round,fill opacity=0.00,] ( 55.20, 55.59) -- ( 55.20,215.35);

\draw[color=drawColor,line cap=round,line join=round,fill opacity=0.00,] ( 55.20, 55.59) -- ( 49.20, 55.59);

\draw[color=drawColor,line cap=round,line join=round,fill opacity=0.00,] ( 55.20, 87.54) -- ( 49.20, 87.54);

\draw[color=drawColor,line cap=round,line join=round,fill opacity=0.00,] ( 55.20,119.50) -- ( 49.20,119.50);

\draw[color=drawColor,line cap=round,line join=round,fill opacity=0.00,] ( 55.20,151.45) -- ( 49.20,151.45);

\draw[color=drawColor,line cap=round,line join=round,fill opacity=0.00,] ( 55.20,183.40) -- ( 49.20,183.40);

\draw[color=drawColor,line cap=round,line join=round,fill opacity=0.00,] ( 55.20,215.35) -- ( 49.20,215.35);

\node[rotate= 90.00,color=drawColor,anchor=base,inner sep=0pt, outer sep=0pt, scale=  0.90] at ( 43.20, 55.59) {0.0%
};

\node[rotate= 90.00,color=drawColor,anchor=base,inner sep=0pt, outer sep=0pt, scale=  0.90] at ( 43.20, 87.54) {0.2%
};

\node[rotate= 90.00,color=drawColor,anchor=base,inner sep=0pt, outer sep=0pt, scale=  0.90] at ( 43.20,119.50) {0.4%
};

\node[rotate= 90.00,color=drawColor,anchor=base,inner sep=0pt, outer sep=0pt, scale=  0.90] at ( 43.20,151.45) {0.6%
};

\node[rotate= 90.00,color=drawColor,anchor=base,inner sep=0pt, outer sep=0pt, scale=  0.90] at ( 43.20,183.40) {0.8%
};

\node[rotate= 90.00,color=drawColor,anchor=base,inner sep=0pt, outer sep=0pt, scale=  0.90] at ( 43.20,215.35) {1.0%
};

\draw[color=drawColor,line cap=round,line join=round,fill opacity=0.00,] ( 55.20, 49.20) --
    (250.54, 49.20) --
    (250.54,221.74) --
    ( 55.20,221.74) --
    ( 55.20, 49.20);
\end{scope}
\begin{scope}
\path[clip] (  0.00,  0.00) rectangle (252.94,252.94);
\definecolor[named]{drawColor}{rgb}{0.60,0.27,0.19}
\definecolor[named]{fillColor}{rgb}{0.88,0.08,0.52}
\definecolor[named]{drawColor}{rgb}{0.00,0.00,0.00}

\node[color=drawColor,anchor=base,inner sep=0pt, outer sep=0pt, scale=  1.00] at (152.87,  1.20) {$\gruen{\nz}$ resp. $\gruen{\nn}$%
};

\node[rotate= 90.00,color=drawColor,anchor=base,inner sep=0pt, outer sep=0pt, scale=  1.00] at ( 19.20,135.47) {$\rot{\yz}$ resp. $\rot{\yn}$%
};
\end{scope}

\draw[-stealth,very thick,gray] (110,175) -- (220,175) node [above,midway,gray] {12 out of 16};
}
\uncover<4->{
\draw[-stealth,very thick] (110,102) -- (220,167) node [above,midway,sloped] {16 out of 16};
}
\end{tikzpicture}
\end{column}
\begin{column}{0.45\textwidth}
\uncover<1->{%
\begin{block}{no conflict:}
prior $\nzg = 8$, $\yzr = 0.75$\\
data $s/n = 12/16 = 0.75$
\end{block}
} %
\uncover<2->{%
\vspace*{-1.5ex}\centerline{\color{lmugreen} $\blacktriangledown$}\vspace*{-1.5ex}
\begin{block}{}
$\nng = 24$, $\ynr = 0.75$
\end{block}
} %
\uncover<4->{%
\vspace*{-1.5ex}\centerline{\color{lmugreen} $\blacktriangle$}\vspace*{-1.5ex}
}
\uncover<3->{%
\begin{block}{\pdc:}
prior $\nzg = 8$, $\yzr = 0.25$\\
data $s/n = 16/16 = 1$
\end{block}
} %
\uncover<0>{%
\vspace*{-1.5ex}\centerline{\color{lmugreen} $\blacktriangledown$}\vspace*{-1.5ex}
\begin{block}{}
$\nng \in [20, 24]$, $\ynr \in [0.73, 0.86]$
\end{block}
} %
%\uncover<5->{ %
%\then\ same predictive prob.\ P!
%}
\end{column}
\end{columns}
}


\subsection{Canonical Conjugate Priors}

\frame{\frametitle{Canonical Conjugate Priors}

\uncover<1->{%
%The multinomial is an example for a \alert{canonical exponential family} %(Bernardo \& Smith, 2000)
Weighted average structure is underneath \alert{all common conjugate priors}
for exponential family sampling distributions! %\\[2ex]

\begin{block}%
{$(x_1, \ldots, x_n) = \x \stackrel{iid}{\sim}$ canonical exponential family} %(Bernardo and Smith, 1994)
%, i.e.
\begin{align*}
p(\x \mid \vartheta) &\propto \exp\big\{\langle \psib, \tau(\x) \rangle - n \bpsib \big\}
\qquad \Big[ \psib \text{ transformation of } \vartheta \Big]
\end{align*}
(includes Binomial, Multinomial, Normal, Poisson, Exponential, \ldots )%\\[2ex]
\end{block}}
\vspace*{-2ex}%
\begin{align*}
\uncover<2->{%
&\text{\play\ conjugate prior:} &
p(\psib\mid\nzg,\yzr) \hspace*{2ex}
&\propto \exp\big\{ \nzg \big[\langle \psib, \yzr \rangle - \bpsib\big]\big\} \\}
\uncover<3->{%
&\text{\play\ (conjugate) posterior:} &
p(\psib\mid\nzg,\yzr,\vec{x})
%p(\psib\mid\nng,\ynr) \hspace*{2ex}
&\propto \exp\big\{ \nng \big[\langle \psib, \ynr \rangle - \bpsib\big]\big\} }%\quad \mbox{where}
\end{align*}
\begin{align*}
\uncover<3->{%
\mbox{where}\quad \ynr &= \frac{\nzg}{\nzg + n} \cdot \yzr + \frac{n}{\nzg + n} \cdot \frac{\tau(\x)}{n}
& &\mbox{and} & \nng &= \nzg + n }
\end{align*}
\note<3>[item]{in the Beta-Binomial model, $\tau(\x) = s$}
}

\frame{\frametitle{Canonical Conjugate Priors}

\vspace*{-2ex}%
\uncover<1->{%
\begin{align*}
&\text{\play\ (conjugate) posterior:} &
%p(\psib\mid\nzg,\yzr,\vec{x})
p(\psib\mid\nng,\ynr) %\hspace*{2ex}
&\propto \exp\big\{ \nng \big[\langle \psib, \ynr \rangle - \bpsib\big]\big\} %\quad \mbox{where}
\end{align*}
\vspace*{-2ex}
\begin{align*}
\mbox{where}\quad \ynr &= \frac{\nzg}{\nzg + n} \cdot \yzr + \frac{n}{\nzg + n} \cdot \frac{\tau(\x)}{n}
& &\mbox{and} & \nng &= \nzg + n
\end{align*}}%
\vspace*{-2ex}%
\uncover<2->{%
\begin{block}{Interpretation}
  \begin{itemize}
  \item $\nzg$ = determines \alert{spread} and \alert{learning speed}
  \item $\yzr$ = \alert{prior expectation} of $\tau(\x)/n$ %$\frac{\tau(\vec{x})}{n}$
  \end{itemize}
\end{block}}%
\uncover<3->{%
\begin{block}{Example: Scaled Normal Data}
\begin{tabular}{r|lcl}
Data :           & $\x\mid\mu$        & $\sim$ & $\norm(\mu,1)$\\[0.5ex] %\sigma_0^2)$ \quad ($\sigma_0^2$ known)\\[0.5ex]
conjugate prior: & $\mu\mid\nzg,\yzr$ & $\sim$ & $\norm(\yzr, 1/\nzg)$  \\[0.5ex]
\cline{1-4}
posterior:       & $\mu\mid\nng,\ynr$ & $\sim$ & $\norm(\ynr, 1/\nng)$ \quad ($\frac{\tau(\x)}{n} = \bar{\x}$)\rule{0ex}{2.5ex} %& \\[0.5ex]
\end{tabular}
\end{block}}
}


\section{Generalised Bayesian Inference}

\subsection{Why Generalise Bayesian Inference?}

\begin{frame}[label=foundational-back]{Why Generalise Bayesian Inference?}
Bayesian theory lacks the ability to specifiy the degree of uncertainty in probability statements
encoded in a (prior, posterior) distribution.\\[1.5ex]

\uncover<2->{%
Variance or stretch of a distribution for describing uncertainty?\\[1.5ex]}
%but this mostly does not work in conjugate updating as just shown.
%but this has mostly no influence on the behaviour to prior-data conflict.

\uncover<3->{%
\then\ \parbox[t]{0.9\textwidth}{Does not work in the case of prior-data conflict:\\
In conjugate updating, the posterior variance does not depend
on the degree of prior-data conflict in most cases.}\\[1.5ex]
\then\ How to express the precision of a probability statement?\\[1.5ex]}

\vfill
{\color{gray}%
\begin{columns}
\begin{column}{0.5\textwidth}
\hyperlink{foundational-app}{\small Foundational arguments regarding over-precision of the classical Bayesian framework}
\end{column}
\begin{column}{0.5\textwidth}
\hyperlink{hierarchical-app}{\small Comments on hierarchical modelling}
\end{column}
\end{columns}
}
\end{frame}


\frame{\frametitle{Imprecision}

%\hspace*{6ex}\parbox[t]{0.9\textwidth}{
\textbf{Add \alert{imprecision} as new model dimension:\\
\alert{Sets of priors} model uncertainty in probability statements\\
and allow to better model partial information on $\vartheta$}\\[2ex]
%}

\uncover<2->{%
\begin{block}{Interpretation}
\centerline{smaller sets $=$ more precise probability statements}
\vspace*{1ex}
\parbox[t]{0.45\textwidth}{\centering \textbf{Lottery A}\\
                          Number of winning tickets:\\
                          exactly known as 5 out of 100\\
                          \then\ $P(\text{win}) = 5/100$}
\qquad
\parbox[t]{0.45\textwidth}{\centering \textbf{Lottery B}\\
                          Number of winning tickets:\\
                          not exactly known, supposedly\\
                          between 1 and 7 out of 100\\
                          \then\ $P(\text{win}) = [1/100,\, 7/100]$}
\end{block}
}
%\uncover<3->{%
%\centerline{shorter intervals \thenthen\ more precise probability statements}
%}
}

\subsection{Sets of Priors}

\frame{\frametitle{Bayesian Inference with Sets of Priors}

\uncover<1->{%
\begin{block}{Standard Bayesian inference procedure}
\centerline{prior + likelihood\hspace*{1ex} \then\ \ posterior}
using Bayes' Rule\\
All inferences are based on the posterior%\\
%(e.g., point estimate = $\E[\psib\mid\nng,\ynr]$)
\end{block}}
\uncover<2->{%
\begin{block}{Generalised Bayesian inference procedure}
\centerline{set of priors + likelihood\hspace*{1ex} \then\ \ set of posteriors}
\emph{Coherence} (consistency of inferences) ensured by using\\
\emph{Generalised Bayes' Rule} (GBR, Walley 1991)\\
= element-wise application of Bayes' Rule\\ %\emph{\small (set of priors must be convex)}\\
%(e.g., point estimate = $\big[\El[\psib\mid\nng,\ynr],\Eu[\psib\mid\nng,\ynr]\big]$)
All inferences are based on the set of posteriors
\end{block}}
\uncover<3->{\centerline{Let hyperparameters $(\nzg, \yzr)$ vary in a set $\PZc$ \then\ \alert{set of priors} $\MZ$}}
}

\begin{frame}{Bayesian Inference with Sets of Priors}

Set of posteriors $\MN$ via $\PNc = \big\{(\nng,\bynr) \colon (\nzg,\byzr) \in \PZc \big\}$
\begin{align*}
\text{single posterior $p(\nng, \ynr)$}
 &\quad\mbox{\then }\quad 
\text{set of posteriors $\MN$ (via $\PNc$)} \\
\E[\psib\mid\nng,\ynr]
 &\quad\mbox{\then }\quad 
\big[\El[\psib\mid\PNc],\Eu[\psib\mid\PNc]\big] \\
P(\psib\in A\mid\nng,\ynr)
 &\quad\mbox{\then }\quad 
\big[\Pl[\psib\in A\mid\PNc],\Pu[\psib\in A\mid\PNc]\big] \\
\text{HPD interval}
 &\quad\mbox{\then }\quad 
\text{union of HPD intervals}
%\E[\theta]  &\quad\mbox{\then }\quad \big[\LE[\theta],\, \UE[\theta]\big] \\
%= \left[ \min_{p \in {\cal M}_\theta} \E_{p}[\theta],\,
%         \max_{p \in {\cal M}_\theta} \E_{p}[\theta]\right]\\
%P(\theta\! \in\! A) &\quad\mbox{\then }\quad \big[\ul{P}(\theta\! \in\! A),\, \ol{P}(\theta\! \in\! A)\big]
%= \big[\! \min P_{\!\!p}(\theta\! \in\! A),\,
%          \max P_{\!\!p}(\theta\! \in\! A)\big]
\end{align*}
\begin{center}
%Lower/upper posterior expectation by min/max over set of posteriors %$\PZc$
Lower/upper bounds by min/max over set of posteriors %$\PZc$
\end{center}
\end{frame}


\begin{frame}[label=commute-back]{Taking the Convex Hull as the Set of Priors}

\uncover<1->{%
\begin{block}{Convex Set of Priors}
%\vspace*{-1ex}
%\begin{align*}
%\MZ &= \text{conv}\big(\big\{p(\psib\mid\nzg,\byzr) \colon (\nzg,\byzr) \in \PZc\big\}\big)\vspace*{-0.5ex}%
%\end{align*}
\centerline{$\MZ = \text{conv}\big(\big\{p(\psib\mid\nzg,\byzr) \colon (\nzg,\byzr) \in \PZc\big\}\big)$}
%Model for prior information =
$\MZ$ = {\small finite convex mixtures of canonical conjugate priors defined by} $\PZc$\\
%Convex set of probability distributions %\\ %{\small(Walley 1991)}\\
%of inference procedure
%(as in standard Bayesian inference procedure)
\end{block}}

\uncover<2->{%
\hyperlink{commute-app}{Updating \& mixture commute} \then\ set of posteriors can be written as\ldots

\begin{block}{Convex Set of Posteriors}
%\vspace*{-1ex}
%\begin{align*}
%\MN &= \text{conv}\big(\big\{p(\psib\mid\nng,\bynr) \colon (\nng,\bynr) \in \PNc\big\}\big)\vspace*{-0.5ex}%
%\end{align*}
\centerline{$\MN = \text{conv}\big(\big\{p(\psib\mid\nng,\bynr) \colon (\nng,\bynr) \in \PNc\big\}\big)$}
%where $\PNc = \big\{(\nng,\bynr) \colon (\nzg,\byzr) \in \PZc \big\}$.\\
%Set of posteriors = 
$\MN$ = finite convex mixtures of canonical conjugate posteriors\\ defined by set of updated hyperparameters $\PNc$
\end{block}}

\uncover<3->{%
Convex sets make the procedure very general (mixture distributions),\\
but are useful only for inferences that are \emph{linear} in the posteriors\\
{\small (expectations: yes, variances: no)}}
\note<3>[item]{so, convexity yes/no is a modelling choice}
\end{frame}


\frame{\frametitle{Imprecise BBM with $\nzg$ fixed: \parbox{0.44\textwidth}{\small IDM (Walley 1996)\\ Quaghebeur \& de Cooman (2005)}}
\hspace*{-12ex}
\begin{columns}%[T]
\begin{column}{0.55\textwidth}
\begin{tikzpicture}
\pgftransformscale{0.025}
\uncover<1>{
% Created by tikzDevice version 0.5.0 on 2011-07-19 12:27:34
\begin{scope}
\path[clip] ( 55.20, 49.20) rectangle (250.54,221.74);
\definecolor[named]{drawColor}{rgb}{0.38,0.64,0.19}
\definecolor[named]{fillColor}{rgb}{0.88,0.08,0.52}
\definecolor[named]{drawColor}{rgb}{0.75,0.75,0.75}
\definecolor[named]{fillColor}{rgb}{0.75,0.75,0.75}

\draw[color=drawColor,line cap=rect,line join=round,fill=fillColor,] ( 96.89, 95.53) circle (  2.25);

\draw[color=drawColor,line cap=rect,line join=round,fill=fillColor,] ( 96.89,175.41) circle (  2.25);

\draw[color=drawColor,line cap=rect,line join=round,fill=fillColor,] (234.70,175.41) circle (  2.25);
\end{scope}
\begin{scope}
\path[clip] (  0.00,  0.00) rectangle (252.94,252.94);
\definecolor[named]{drawColor}{rgb}{0.38,0.64,0.19}
\definecolor[named]{fillColor}{rgb}{0.88,0.08,0.52}
\definecolor[named]{drawColor}{rgb}{0.00,0.00,0.00}

\draw[color=drawColor,line cap=rect,line join=round,fill opacity=0.00,] ( 71.05, 49.20) -- (243.31, 49.20);

\draw[color=drawColor,line cap=rect,line join=round,fill opacity=0.00,] ( 71.05, 49.20) -- ( 71.05, 43.20);

\draw[color=drawColor,line cap=rect,line join=round,fill opacity=0.00,] (114.11, 49.20) -- (114.11, 43.20);

\draw[color=drawColor,line cap=rect,line join=round,fill opacity=0.00,] (157.18, 49.20) -- (157.18, 43.20);

\draw[color=drawColor,line cap=rect,line join=round,fill opacity=0.00,] (200.24, 49.20) -- (200.24, 43.20);

\draw[color=drawColor,line cap=rect,line join=round,fill opacity=0.00,] (243.31, 49.20) -- (243.31, 43.20);

\node[color=drawColor,anchor=base,inner sep=0pt, outer sep=0pt, scale=  0.90] at ( 71.05, 25.20) {5%
};

\node[color=drawColor,anchor=base,inner sep=0pt, outer sep=0pt, scale=  0.90] at (114.11, 25.20) {10%
};

\node[color=drawColor,anchor=base,inner sep=0pt, outer sep=0pt, scale=  0.90] at (157.18, 25.20) {15%
};

\node[color=drawColor,anchor=base,inner sep=0pt, outer sep=0pt, scale=  0.90] at (200.24, 25.20) {20%
};

\node[color=drawColor,anchor=base,inner sep=0pt, outer sep=0pt, scale=  0.90] at (243.31, 25.20) {25%
};

\draw[color=drawColor,line cap=rect,line join=round,fill opacity=0.00,] ( 55.20, 55.59) -- ( 55.20,215.35);

\draw[color=drawColor,line cap=rect,line join=round,fill opacity=0.00,] ( 55.20, 55.59) -- ( 49.20, 55.59);

\draw[color=drawColor,line cap=rect,line join=round,fill opacity=0.00,] ( 55.20, 87.54) -- ( 49.20, 87.54);

\draw[color=drawColor,line cap=rect,line join=round,fill opacity=0.00,] ( 55.20,119.50) -- ( 49.20,119.50);

\draw[color=drawColor,line cap=rect,line join=round,fill opacity=0.00,] ( 55.20,151.45) -- ( 49.20,151.45);

\draw[color=drawColor,line cap=rect,line join=round,fill opacity=0.00,] ( 55.20,183.40) -- ( 49.20,183.40);

\draw[color=drawColor,line cap=rect,line join=round,fill opacity=0.00,] ( 55.20,215.35) -- ( 49.20,215.35);

\node[rotate= 90.00,color=drawColor,anchor=base,inner sep=0pt, outer sep=0pt, scale=  0.90] at ( 43.20, 55.59) {0.0%
};

\node[rotate= 90.00,color=drawColor,anchor=base,inner sep=0pt, outer sep=0pt, scale=  0.90] at ( 43.20, 87.54) {0.2%
};

\node[rotate= 90.00,color=drawColor,anchor=base,inner sep=0pt, outer sep=0pt, scale=  0.90] at ( 43.20,119.50) {0.4%
};

\node[rotate= 90.00,color=drawColor,anchor=base,inner sep=0pt, outer sep=0pt, scale=  0.90] at ( 43.20,151.45) {0.6%
};

\node[rotate= 90.00,color=drawColor,anchor=base,inner sep=0pt, outer sep=0pt, scale=  0.90] at ( 43.20,183.40) {0.8%
};

\node[rotate= 90.00,color=drawColor,anchor=base,inner sep=0pt, outer sep=0pt, scale=  0.90] at ( 43.20,215.35) {1.0%
};

\draw[color=drawColor,line cap=rect,line join=round,fill opacity=0.00,] ( 55.20, 49.20) --
    (250.54, 49.20) --
    (250.54,221.74) --
    ( 55.20,221.74) --
    ( 55.20, 49.20);
\end{scope}
\begin{scope}
\path[clip] (  0.00,  0.00) rectangle (252.94,252.94);
\definecolor[named]{drawColor}{rgb}{0.38,0.64,0.19}
\definecolor[named]{fillColor}{rgb}{0.88,0.08,0.52}
\definecolor[named]{drawColor}{rgb}{0.00,0.00,0.00}

\node[color=drawColor,anchor=base,inner sep=0pt, outer sep=0pt, scale=  1.00] at (152.87,  1.20) {$\gruen{\nz}$ resp. $\gruen{\nn}$%
};

\node[rotate= 90.00,color=drawColor,anchor=base,inner sep=0pt, outer sep=0pt, scale=  1.00] at ( 19.20,135.47) {$\rot{\yz}$ resp. $\rot{\yn}$%
};
\end{scope}
\begin{scope}
\path[clip] ( 55.20, 49.20) rectangle (250.54,221.74);
\definecolor[named]{drawColor}{rgb}{0.38,0.64,0.19}
\definecolor[named]{fillColor}{rgb}{0.88,0.08,0.52}
\definecolor[named]{drawColor}{rgb}{0.00,0.00,0.00}
\definecolor[named]{fillColor}{rgb}{0.75,0.75,0.75}

\draw[color=drawColor,line width= 1.2pt,line cap=rect,line join=round,fill=fillColor,] ( 96.89,167.43) --
    ( 96.89,167.43) --
    ( 96.89,183.40) --
    ( 96.89,183.40) --
    cycle;
\end{scope}

}
\uncover<2>{
% Created by tikzDevice version 0.5.0 on 2011-07-19 12:27:41
\begin{scope}
\path[clip] ( 55.20, 49.20) rectangle (250.54,221.74);
\definecolor[named]{drawColor}{rgb}{0.38,0.14,0.19}
\definecolor[named]{fillColor}{rgb}{0.88,0.08,0.52}
\definecolor[named]{drawColor}{rgb}{0.75,0.75,0.75}
\definecolor[named]{fillColor}{rgb}{0.75,0.75,0.75}

\draw[color=drawColor,line cap=rect,line join=round,fill=fillColor,] ( 96.89, 95.53) circle (  2.25);

\draw[color=drawColor,line cap=rect,line join=round,fill=fillColor,] ( 96.89,175.41) circle (  2.25);

\draw[color=drawColor,line cap=rect,line join=round,fill=fillColor,] (234.70,175.41) circle (  2.25);
\end{scope}
\begin{scope}
\path[clip] (  0.00,  0.00) rectangle (252.94,252.94);
\definecolor[named]{drawColor}{rgb}{0.38,0.14,0.19}
\definecolor[named]{fillColor}{rgb}{0.88,0.08,0.52}
\definecolor[named]{drawColor}{rgb}{0.00,0.00,0.00}

\draw[color=drawColor,line cap=rect,line join=round,fill opacity=0.00,] ( 71.05, 49.20) -- (243.31, 49.20);

\draw[color=drawColor,line cap=rect,line join=round,fill opacity=0.00,] ( 71.05, 49.20) -- ( 71.05, 43.20);

\draw[color=drawColor,line cap=rect,line join=round,fill opacity=0.00,] (114.11, 49.20) -- (114.11, 43.20);

\draw[color=drawColor,line cap=rect,line join=round,fill opacity=0.00,] (157.18, 49.20) -- (157.18, 43.20);

\draw[color=drawColor,line cap=rect,line join=round,fill opacity=0.00,] (200.24, 49.20) -- (200.24, 43.20);

\draw[color=drawColor,line cap=rect,line join=round,fill opacity=0.00,] (243.31, 49.20) -- (243.31, 43.20);

\node[color=drawColor,anchor=base,inner sep=0pt, outer sep=0pt, scale=  0.90] at ( 71.05, 25.20) {5%
};

\node[color=drawColor,anchor=base,inner sep=0pt, outer sep=0pt, scale=  0.90] at (114.11, 25.20) {10%
};

\node[color=drawColor,anchor=base,inner sep=0pt, outer sep=0pt, scale=  0.90] at (157.18, 25.20) {15%
};

\node[color=drawColor,anchor=base,inner sep=0pt, outer sep=0pt, scale=  0.90] at (200.24, 25.20) {20%
};

\node[color=drawColor,anchor=base,inner sep=0pt, outer sep=0pt, scale=  0.90] at (243.31, 25.20) {25%
};

\draw[color=drawColor,line cap=rect,line join=round,fill opacity=0.00,] ( 55.20, 55.59) -- ( 55.20,215.35);

\draw[color=drawColor,line cap=rect,line join=round,fill opacity=0.00,] ( 55.20, 55.59) -- ( 49.20, 55.59);

\draw[color=drawColor,line cap=rect,line join=round,fill opacity=0.00,] ( 55.20, 87.54) -- ( 49.20, 87.54);

\draw[color=drawColor,line cap=rect,line join=round,fill opacity=0.00,] ( 55.20,119.50) -- ( 49.20,119.50);

\draw[color=drawColor,line cap=rect,line join=round,fill opacity=0.00,] ( 55.20,151.45) -- ( 49.20,151.45);

\draw[color=drawColor,line cap=rect,line join=round,fill opacity=0.00,] ( 55.20,183.40) -- ( 49.20,183.40);

\draw[color=drawColor,line cap=rect,line join=round,fill opacity=0.00,] ( 55.20,215.35) -- ( 49.20,215.35);

\node[rotate= 90.00,color=drawColor,anchor=base,inner sep=0pt, outer sep=0pt, scale=  0.90] at ( 43.20, 55.59) {0.0%
};

\node[rotate= 90.00,color=drawColor,anchor=base,inner sep=0pt, outer sep=0pt, scale=  0.90] at ( 43.20, 87.54) {0.2%
};

\node[rotate= 90.00,color=drawColor,anchor=base,inner sep=0pt, outer sep=0pt, scale=  0.90] at ( 43.20,119.50) {0.4%
};

\node[rotate= 90.00,color=drawColor,anchor=base,inner sep=0pt, outer sep=0pt, scale=  0.90] at ( 43.20,151.45) {0.6%
};

\node[rotate= 90.00,color=drawColor,anchor=base,inner sep=0pt, outer sep=0pt, scale=  0.90] at ( 43.20,183.40) {0.8%
};

\node[rotate= 90.00,color=drawColor,anchor=base,inner sep=0pt, outer sep=0pt, scale=  0.90] at ( 43.20,215.35) {1.0%
};

\draw[color=drawColor,line cap=rect,line join=round,fill opacity=0.00,] ( 55.20, 49.20) --
    (250.54, 49.20) --
    (250.54,221.74) --
    ( 55.20,221.74) --
    ( 55.20, 49.20);
\end{scope}
\begin{scope}
\path[clip] (  0.00,  0.00) rectangle (252.94,252.94);
\definecolor[named]{drawColor}{rgb}{0.38,0.14,0.19}
\definecolor[named]{fillColor}{rgb}{0.88,0.08,0.52}
\definecolor[named]{drawColor}{rgb}{0.00,0.00,0.00}

\node[color=drawColor,anchor=base,inner sep=0pt, outer sep=0pt, scale=  1.00] at (152.87,  1.20) {$\gruen{\nz}$ resp. $\gruen{\nn}$%
};

\node[rotate= 90.00,color=drawColor,anchor=base,inner sep=0pt, outer sep=0pt, scale=  1.00] at ( 19.20,135.47) {$\rot{\yz}$ resp. $\rot{\yn}$%
};
\end{scope}
\begin{scope}
\path[clip] ( 55.20, 49.20) rectangle (250.54,221.74);
\definecolor[named]{drawColor}{rgb}{0.38,0.14,0.19}
\definecolor[named]{fillColor}{rgb}{0.88,0.08,0.52}
\definecolor[named]{drawColor}{rgb}{0.00,0.00,0.00}
\definecolor[named]{fillColor}{rgb}{0.75,0.75,0.75}

\draw[color=drawColor,line width= 1.2pt,line cap=rect,line join=round,fill=fillColor,] ( 96.89,167.43) --
    ( 96.89,167.43) --
    ( 96.89,183.40) --
    ( 96.89,183.40) --
    cycle;

\draw[color=drawColor,line width= 1.2pt,line cap=rect,line join=round,fill=fillColor,] (234.70,172.75) --
    (234.70,172.75) --
    (234.70,172.75) --
    (234.70,172.75) --
    (234.70,172.75) --
    (234.70,172.75) --
    (234.70,172.75) --
    (234.70,172.75) --
    (234.70,172.75) --
    (234.70,172.75) --
    (234.70,172.75) --
    (234.70,172.75) --
    (234.70,172.75) --
    (234.70,172.75) --
    (234.70,172.75) --
    (234.70,172.75) --
    (234.70,172.75) --
    (234.70,172.75) --
    (234.70,172.75) --
    (234.70,172.75) --
    (234.70,172.75) --
    (234.70,172.75) --
    (234.70,172.75) --
    (234.70,172.75) --
    (234.70,172.75) --
    (234.70,172.75) --
    (234.70,172.75) --
    (234.70,172.75) --
    (234.70,172.75) --
    (234.70,172.75) --
    (234.70,172.75) --
    (234.70,172.75) --
    (234.70,172.75) --
    (234.70,172.75) --
    (234.70,172.75) --
    (234.70,172.75) --
    (234.70,172.75) --
    (234.70,172.75) --
    (234.70,172.75) --
    (234.70,172.75) --
    (234.70,172.75) --
    (234.70,172.75) --
    (234.70,172.75) --
    (234.70,172.75) --
    (234.70,172.75) --
    (234.70,172.75) --
    (234.70,172.75) --
    (234.70,172.75) --
    (234.70,172.75) --
    (234.70,172.75) --
    (234.70,172.75) --
    (234.70,172.75) --
    (234.70,172.75) --
    (234.70,172.75) --
    (234.70,172.75) --
    (234.70,172.75) --
    (234.70,172.75) --
    (234.70,172.75) --
    (234.70,172.75) --
    (234.70,172.75) --
    (234.70,172.75) --
    (234.70,172.75) --
    (234.70,172.75) --
    (234.70,172.75) --
    (234.70,172.75) --
    (234.70,172.75) --
    (234.70,172.75) --
    (234.70,172.75) --
    (234.70,172.75) --
    (234.70,172.75) --
    (234.70,172.75) --
    (234.70,172.75) --
    (234.70,172.75) --
    (234.70,172.75) --
    (234.70,172.75) --
    (234.70,172.75) --
    (234.70,172.75) --
    (234.70,172.75) --
    (234.70,172.75) --
    (234.70,172.75) --
    (234.70,172.75) --
    (234.70,172.75) --
    (234.70,172.75) --
    (234.70,172.75) --
    (234.70,172.75) --
    (234.70,172.75) --
    (234.70,172.75) --
    (234.70,172.75) --
    (234.70,172.75) --
    (234.70,172.75) --
    (234.70,172.75) --
    (234.70,172.75) --
    (234.70,172.75) --
    (234.70,172.75) --
    (234.70,172.75) --
    (234.70,172.75) --
    (234.70,172.75) --
    (234.70,172.75) --
    (234.70,172.75) --
    (234.70,172.75) --
    (234.70,178.08) --
    (234.70,178.08) --
    (234.70,178.08) --
    (234.70,178.08) --
    (234.70,178.08) --
    (234.70,178.08) --
    (234.70,178.08) --
    (234.70,178.08) --
    (234.70,178.08) --
    (234.70,178.08) --
    (234.70,178.08) --
    (234.70,178.08) --
    (234.70,178.08) --
    (234.70,178.08) --
    (234.70,178.08) --
    (234.70,178.08) --
    (234.70,178.08) --
    (234.70,178.08) --
    (234.70,178.08) --
    (234.70,178.08) --
    (234.70,178.08) --
    (234.70,178.08) --
    (234.70,178.08) --
    (234.70,178.08) --
    (234.70,178.08) --
    (234.70,178.08) --
    (234.70,178.08) --
    (234.70,178.08) --
    (234.70,178.08) --
    (234.70,178.08) --
    (234.70,178.08) --
    (234.70,178.08) --
    (234.70,178.08) --
    (234.70,178.08) --
    (234.70,178.08) --
    (234.70,178.08) --
    (234.70,178.08) --
    (234.70,178.08) --
    (234.70,178.08) --
    (234.70,178.08) --
    (234.70,178.08) --
    (234.70,178.08) --
    (234.70,178.08) --
    (234.70,178.08) --
    (234.70,178.08) --
    (234.70,178.08) --
    (234.70,178.08) --
    (234.70,178.08) --
    (234.70,178.08) --
    (234.70,178.08) --
    (234.70,178.08) --
    (234.70,178.08) --
    (234.70,178.08) --
    (234.70,178.08) --
    (234.70,178.08) --
    (234.70,178.08) --
    (234.70,178.08) --
    (234.70,178.08) --
    (234.70,178.08) --
    (234.70,178.08) --
    (234.70,178.08) --
    (234.70,178.08) --
    (234.70,178.08) --
    (234.70,178.08) --
    (234.70,178.08) --
    (234.70,178.08) --
    (234.70,178.08) --
    (234.70,178.08) --
    (234.70,178.08) --
    (234.70,178.08) --
    (234.70,178.08) --
    (234.70,178.08) --
    (234.70,178.08) --
    (234.70,178.08) --
    (234.70,178.08) --
    (234.70,178.08) --
    (234.70,178.08) --
    (234.70,178.08) --
    (234.70,178.08) --
    (234.70,178.08) --
    (234.70,178.08) --
    (234.70,178.08) --
    (234.70,178.08) --
    (234.70,178.08) --
    (234.70,178.08) --
    (234.70,178.08) --
    (234.70,178.08) --
    (234.70,178.08) --
    (234.70,178.08) --
    (234.70,178.08) --
    (234.70,178.08) --
    (234.70,178.08) --
    (234.70,178.08) --
    (234.70,178.08) --
    (234.70,178.08) --
    (234.70,178.08) --
    (234.70,178.08) --
    (234.70,178.08) --
    (234.70,178.08) --
    (234.70,178.08) --
    cycle;
\end{scope}

\draw[-stealth,very thick] (110,175) -- (220,175) node [above,midway] {12 out of 16};
}
\uncover<3->{
% Created by tikzDevice version 0.5.0 on 2011-07-19 12:27:48
\begin{scope}
\path[clip] ( 55.20, 49.20) rectangle (250.54,221.74);
\definecolor[named]{drawColor}{rgb}{0.19,0.42,0.19}
\definecolor[named]{fillColor}{rgb}{0.88,0.08,0.52}
\definecolor[named]{drawColor}{rgb}{0.75,0.75,0.75}
\definecolor[named]{fillColor}{rgb}{0.75,0.75,0.75}

\draw[color=drawColor,line cap=rect,line join=round,fill=fillColor,] ( 96.89, 95.53) circle (  2.25);

\draw[color=drawColor,line cap=rect,line join=round,fill=fillColor,] ( 96.89,175.41) circle (  2.25);

\draw[color=drawColor,line cap=rect,line join=round,fill=fillColor,] (234.70,175.41) circle (  2.25);
\end{scope}
\begin{scope}
\path[clip] (  0.00,  0.00) rectangle (252.94,252.94);
\definecolor[named]{drawColor}{rgb}{0.19,0.42,0.19}
\definecolor[named]{fillColor}{rgb}{0.88,0.08,0.52}
\definecolor[named]{drawColor}{rgb}{0.00,0.00,0.00}

\draw[color=drawColor,line cap=rect,line join=round,fill opacity=0.00,] ( 71.05, 49.20) -- (243.31, 49.20);

\draw[color=drawColor,line cap=rect,line join=round,fill opacity=0.00,] ( 71.05, 49.20) -- ( 71.05, 43.20);

\draw[color=drawColor,line cap=rect,line join=round,fill opacity=0.00,] (114.11, 49.20) -- (114.11, 43.20);

\draw[color=drawColor,line cap=rect,line join=round,fill opacity=0.00,] (157.18, 49.20) -- (157.18, 43.20);

\draw[color=drawColor,line cap=rect,line join=round,fill opacity=0.00,] (200.24, 49.20) -- (200.24, 43.20);

\draw[color=drawColor,line cap=rect,line join=round,fill opacity=0.00,] (243.31, 49.20) -- (243.31, 43.20);

\node[color=drawColor,anchor=base,inner sep=0pt, outer sep=0pt, scale=  0.90] at ( 71.05, 25.20) {5%
};

\node[color=drawColor,anchor=base,inner sep=0pt, outer sep=0pt, scale=  0.90] at (114.11, 25.20) {10%
};

\node[color=drawColor,anchor=base,inner sep=0pt, outer sep=0pt, scale=  0.90] at (157.18, 25.20) {15%
};

\node[color=drawColor,anchor=base,inner sep=0pt, outer sep=0pt, scale=  0.90] at (200.24, 25.20) {20%
};

\node[color=drawColor,anchor=base,inner sep=0pt, outer sep=0pt, scale=  0.90] at (243.31, 25.20) {25%
};

\draw[color=drawColor,line cap=rect,line join=round,fill opacity=0.00,] ( 55.20, 55.59) -- ( 55.20,215.35);

\draw[color=drawColor,line cap=rect,line join=round,fill opacity=0.00,] ( 55.20, 55.59) -- ( 49.20, 55.59);

\draw[color=drawColor,line cap=rect,line join=round,fill opacity=0.00,] ( 55.20, 87.54) -- ( 49.20, 87.54);

\draw[color=drawColor,line cap=rect,line join=round,fill opacity=0.00,] ( 55.20,119.50) -- ( 49.20,119.50);

\draw[color=drawColor,line cap=rect,line join=round,fill opacity=0.00,] ( 55.20,151.45) -- ( 49.20,151.45);

\draw[color=drawColor,line cap=rect,line join=round,fill opacity=0.00,] ( 55.20,183.40) -- ( 49.20,183.40);

\draw[color=drawColor,line cap=rect,line join=round,fill opacity=0.00,] ( 55.20,215.35) -- ( 49.20,215.35);

\node[rotate= 90.00,color=drawColor,anchor=base,inner sep=0pt, outer sep=0pt, scale=  0.90] at ( 43.20, 55.59) {0.0%
};

\node[rotate= 90.00,color=drawColor,anchor=base,inner sep=0pt, outer sep=0pt, scale=  0.90] at ( 43.20, 87.54) {0.2%
};

\node[rotate= 90.00,color=drawColor,anchor=base,inner sep=0pt, outer sep=0pt, scale=  0.90] at ( 43.20,119.50) {0.4%
};

\node[rotate= 90.00,color=drawColor,anchor=base,inner sep=0pt, outer sep=0pt, scale=  0.90] at ( 43.20,151.45) {0.6%
};

\node[rotate= 90.00,color=drawColor,anchor=base,inner sep=0pt, outer sep=0pt, scale=  0.90] at ( 43.20,183.40) {0.8%
};

\node[rotate= 90.00,color=drawColor,anchor=base,inner sep=0pt, outer sep=0pt, scale=  0.90] at ( 43.20,215.35) {1.0%
};

\draw[color=drawColor,line cap=rect,line join=round,fill opacity=0.00,] ( 55.20, 49.20) --
    (250.54, 49.20) --
    (250.54,221.74) --
    ( 55.20,221.74) --
    ( 55.20, 49.20);
\end{scope}
\begin{scope}
\path[clip] (  0.00,  0.00) rectangle (252.94,252.94);
\definecolor[named]{drawColor}{rgb}{0.19,0.42,0.19}
\definecolor[named]{fillColor}{rgb}{0.88,0.08,0.52}
\definecolor[named]{drawColor}{rgb}{0.00,0.00,0.00}

\node[color=drawColor,anchor=base,inner sep=0pt, outer sep=0pt, scale=  1.00] at (152.87,  1.20) {$\gruen{\nz}$ resp. $\gruen{\nn}$%
};

\node[rotate= 90.00,color=drawColor,anchor=base,inner sep=0pt, outer sep=0pt, scale=  1.00] at ( 19.20,135.47) {$\rot{\yz}$ resp. $\rot{\yn}$%
};
\end{scope}
\begin{scope}
\path[clip] ( 55.20, 49.20) rectangle (250.54,221.74);
\definecolor[named]{drawColor}{rgb}{0.19,0.42,0.19}
\definecolor[named]{fillColor}{rgb}{0.88,0.08,0.52}
\definecolor[named]{drawColor}{rgb}{0.00,0.00,0.00}
\definecolor[named]{fillColor}{rgb}{0.75,0.75,0.75}

\draw[color=drawColor,line width= 1.2pt,line cap=rect,line join=round,fill=fillColor,] ( 96.89,167.43) --
    ( 96.89,167.43) --
    ( 96.89,183.40) --
    ( 96.89,183.40) --
    cycle;

\draw[color=drawColor,line width= 1.2pt,line cap=rect,line join=round,fill=fillColor,] (234.70,172.75) --
    (234.70,172.75) --
    (234.70,172.75) --
    (234.70,172.75) --
    (234.70,172.75) --
    (234.70,172.75) --
    (234.70,172.75) --
    (234.70,172.75) --
    (234.70,172.75) --
    (234.70,172.75) --
    (234.70,172.75) --
    (234.70,172.75) --
    (234.70,172.75) --
    (234.70,172.75) --
    (234.70,172.75) --
    (234.70,172.75) --
    (234.70,172.75) --
    (234.70,172.75) --
    (234.70,172.75) --
    (234.70,172.75) --
    (234.70,172.75) --
    (234.70,172.75) --
    (234.70,172.75) --
    (234.70,172.75) --
    (234.70,172.75) --
    (234.70,172.75) --
    (234.70,172.75) --
    (234.70,172.75) --
    (234.70,172.75) --
    (234.70,172.75) --
    (234.70,172.75) --
    (234.70,172.75) --
    (234.70,172.75) --
    (234.70,172.75) --
    (234.70,172.75) --
    (234.70,172.75) --
    (234.70,172.75) --
    (234.70,172.75) --
    (234.70,172.75) --
    (234.70,172.75) --
    (234.70,172.75) --
    (234.70,172.75) --
    (234.70,172.75) --
    (234.70,172.75) --
    (234.70,172.75) --
    (234.70,172.75) --
    (234.70,172.75) --
    (234.70,172.75) --
    (234.70,172.75) --
    (234.70,172.75) --
    (234.70,172.75) --
    (234.70,172.75) --
    (234.70,172.75) --
    (234.70,172.75) --
    (234.70,172.75) --
    (234.70,172.75) --
    (234.70,172.75) --
    (234.70,172.75) --
    (234.70,172.75) --
    (234.70,172.75) --
    (234.70,172.75) --
    (234.70,172.75) --
    (234.70,172.75) --
    (234.70,172.75) --
    (234.70,172.75) --
    (234.70,172.75) --
    (234.70,172.75) --
    (234.70,172.75) --
    (234.70,172.75) --
    (234.70,172.75) --
    (234.70,172.75) --
    (234.70,172.75) --
    (234.70,172.75) --
    (234.70,172.75) --
    (234.70,172.75) --
    (234.70,172.75) --
    (234.70,172.75) --
    (234.70,172.75) --
    (234.70,172.75) --
    (234.70,172.75) --
    (234.70,172.75) --
    (234.70,172.75) --
    (234.70,172.75) --
    (234.70,172.75) --
    (234.70,172.75) --
    (234.70,172.75) --
    (234.70,172.75) --
    (234.70,172.75) --
    (234.70,172.75) --
    (234.70,172.75) --
    (234.70,172.75) --
    (234.70,172.75) --
    (234.70,172.75) --
    (234.70,172.75) --
    (234.70,172.75) --
    (234.70,172.75) --
    (234.70,172.75) --
    (234.70,172.75) --
    (234.70,172.75) --
    (234.70,172.75) --
    (234.70,178.08) --
    (234.70,178.08) --
    (234.70,178.08) --
    (234.70,178.08) --
    (234.70,178.08) --
    (234.70,178.08) --
    (234.70,178.08) --
    (234.70,178.08) --
    (234.70,178.08) --
    (234.70,178.08) --
    (234.70,178.08) --
    (234.70,178.08) --
    (234.70,178.08) --
    (234.70,178.08) --
    (234.70,178.08) --
    (234.70,178.08) --
    (234.70,178.08) --
    (234.70,178.08) --
    (234.70,178.08) --
    (234.70,178.08) --
    (234.70,178.08) --
    (234.70,178.08) --
    (234.70,178.08) --
    (234.70,178.08) --
    (234.70,178.08) --
    (234.70,178.08) --
    (234.70,178.08) --
    (234.70,178.08) --
    (234.70,178.08) --
    (234.70,178.08) --
    (234.70,178.08) --
    (234.70,178.08) --
    (234.70,178.08) --
    (234.70,178.08) --
    (234.70,178.08) --
    (234.70,178.08) --
    (234.70,178.08) --
    (234.70,178.08) --
    (234.70,178.08) --
    (234.70,178.08) --
    (234.70,178.08) --
    (234.70,178.08) --
    (234.70,178.08) --
    (234.70,178.08) --
    (234.70,178.08) --
    (234.70,178.08) --
    (234.70,178.08) --
    (234.70,178.08) --
    (234.70,178.08) --
    (234.70,178.08) --
    (234.70,178.08) --
    (234.70,178.08) --
    (234.70,178.08) --
    (234.70,178.08) --
    (234.70,178.08) --
    (234.70,178.08) --
    (234.70,178.08) --
    (234.70,178.08) --
    (234.70,178.08) --
    (234.70,178.08) --
    (234.70,178.08) --
    (234.70,178.08) --
    (234.70,178.08) --
    (234.70,178.08) --
    (234.70,178.08) --
    (234.70,178.08) --
    (234.70,178.08) --
    (234.70,178.08) --
    (234.70,178.08) --
    (234.70,178.08) --
    (234.70,178.08) --
    (234.70,178.08) --
    (234.70,178.08) --
    (234.70,178.08) --
    (234.70,178.08) --
    (234.70,178.08) --
    (234.70,178.08) --
    (234.70,178.08) --
    (234.70,178.08) --
    (234.70,178.08) --
    (234.70,178.08) --
    (234.70,178.08) --
    (234.70,178.08) --
    (234.70,178.08) --
    (234.70,178.08) --
    (234.70,178.08) --
    (234.70,178.08) --
    (234.70,178.08) --
    (234.70,178.08) --
    (234.70,178.08) --
    (234.70,178.08) --
    (234.70,178.08) --
    (234.70,178.08) --
    (234.70,178.08) --
    (234.70,178.08) --
    (234.70,178.08) --
    (234.70,178.08) --
    (234.70,178.08) --
    (234.70,178.08) --
    (234.70,178.08) --
    cycle;

\draw[color=drawColor,line width= 1.2pt,line cap=rect,line join=round,fill=fillColor,] ( 96.89, 87.54) --
    ( 96.89, 87.54) --
    ( 96.89,103.52) --
    ( 96.89,103.52) --
    cycle;
\end{scope}

\draw[-stealth,very thick,gray] (110,175) -- (220,175) node [above,midway,gray] {12 out of 16};
}
\uncover<4->{
\draw[-stealth,very thick] (110,102) -- (220,167) node [above,midway,sloped] {16 out of 16};
}
\end{tikzpicture}
\end{column}
\begin{column}{0.45\textwidth}
\uncover<1->{%
\begin{block}{no conflict:}
prior $\nzg = 8$, $\yzr \in [0.7, 0.8]$\\
data $s/n = 12/16 = 0.75$
\end{block}
} %
\uncover<2->{%
\vspace*{-1.5ex}\centerline{\color{lmugreen} $\blacktriangledown$}\vspace*{-1.5ex}
\begin{block}{}
$\nng = 24$, $\ynr \in [0.73, 0.77]$
\end{block}
} %
\uncover<4->{%
\vspace*{-1.5ex}\centerline{\color{lmugreen} $\blacktriangle$}\vspace*{-1.5ex}
}
\uncover<3->{%
\begin{block}{\pdc:}
prior $\nzg = 8$, $\yzr \in [0.2, 0.3]$\\
data $s/n = 16/16 = 1$
\end{block}
} %
\uncover<0>{%
\vspace*{-1.5ex}\centerline{\color{lmugreen} $\blacktriangledown$}\vspace*{-1.5ex}
\begin{block}{}
$\nng \in [20, 24]$, $\ynr \in [0.73, 0.86]$
\end{block}
} %
%\uncover<5->{ %
%\then\ same imprecise predictive prob.\ $[\Pl,\Pu]$!
%}
\end{column}
\end{columns}

}


\frame{\frametitle{Imprecise BBM with $[\nzlg,\nzug]$: \parbox{0.43\textwidth}{\small Walley (1991, \S 5.4.3)\\ Walter \& Augustin (2009)}}
\hspace*{-12ex}
\begin{columns}%[T]
\begin{column}{0.55\textwidth}
\begin{tikzpicture}
\pgftransformscale{0.025}
\uncover<1>{
% Created by tikzDevice version 0.5.0 on 2011-07-19 12:36:59
\begin{scope}
\path[clip] ( 55.20, 49.20) rectangle (250.54,221.74);
\definecolor[named]{drawColor}{rgb}{0.97,1.00,0.19}
\definecolor[named]{fillColor}{rgb}{0.88,0.08,0.52}
\definecolor[named]{drawColor}{rgb}{0.75,0.75,0.75}
\definecolor[named]{fillColor}{rgb}{0.75,0.75,0.75}

\draw[color=drawColor,line cap=round,line join=round,fill=fillColor,] ( 96.89, 95.53) circle (  2.25);

\draw[color=drawColor,line cap=round,line join=round,fill=fillColor,] ( 96.89,175.41) circle (  2.25);

\draw[color=drawColor,line cap=round,line join=round,fill=fillColor,] (234.70,175.41) circle (  2.25);
\end{scope}
\begin{scope}
\path[clip] (  0.00,  0.00) rectangle (252.94,252.94);
\definecolor[named]{drawColor}{rgb}{0.97,1.00,0.19}
\definecolor[named]{fillColor}{rgb}{0.88,0.08,0.52}
\definecolor[named]{drawColor}{rgb}{0.00,0.00,0.00}

\draw[color=drawColor,line cap=round,line join=round,fill opacity=0.00,] ( 71.05, 49.20) -- (243.31, 49.20);

\draw[color=drawColor,line cap=round,line join=round,fill opacity=0.00,] ( 71.05, 49.20) -- ( 71.05, 43.20);

\draw[color=drawColor,line cap=round,line join=round,fill opacity=0.00,] (114.11, 49.20) -- (114.11, 43.20);

\draw[color=drawColor,line cap=round,line join=round,fill opacity=0.00,] (157.18, 49.20) -- (157.18, 43.20);

\draw[color=drawColor,line cap=round,line join=round,fill opacity=0.00,] (200.24, 49.20) -- (200.24, 43.20);

\draw[color=drawColor,line cap=round,line join=round,fill opacity=0.00,] (243.31, 49.20) -- (243.31, 43.20);

\node[color=drawColor,anchor=base,inner sep=0pt, outer sep=0pt, scale=  0.90] at ( 71.05, 25.20) {5%
};

\node[color=drawColor,anchor=base,inner sep=0pt, outer sep=0pt, scale=  0.90] at (114.11, 25.20) {10%
};

\node[color=drawColor,anchor=base,inner sep=0pt, outer sep=0pt, scale=  0.90] at (157.18, 25.20) {15%
};

\node[color=drawColor,anchor=base,inner sep=0pt, outer sep=0pt, scale=  0.90] at (200.24, 25.20) {20%
};

\node[color=drawColor,anchor=base,inner sep=0pt, outer sep=0pt, scale=  0.90] at (243.31, 25.20) {25%
};

\draw[color=drawColor,line cap=round,line join=round,fill opacity=0.00,] ( 55.20, 55.59) -- ( 55.20,215.35);

\draw[color=drawColor,line cap=round,line join=round,fill opacity=0.00,] ( 55.20, 55.59) -- ( 49.20, 55.59);

\draw[color=drawColor,line cap=round,line join=round,fill opacity=0.00,] ( 55.20, 87.54) -- ( 49.20, 87.54);

\draw[color=drawColor,line cap=round,line join=round,fill opacity=0.00,] ( 55.20,119.50) -- ( 49.20,119.50);

\draw[color=drawColor,line cap=round,line join=round,fill opacity=0.00,] ( 55.20,151.45) -- ( 49.20,151.45);

\draw[color=drawColor,line cap=round,line join=round,fill opacity=0.00,] ( 55.20,183.40) -- ( 49.20,183.40);

\draw[color=drawColor,line cap=round,line join=round,fill opacity=0.00,] ( 55.20,215.35) -- ( 49.20,215.35);

\node[rotate= 90.00,color=drawColor,anchor=base,inner sep=0pt, outer sep=0pt, scale=  0.90] at ( 43.20, 55.59) {0.0%
};

\node[rotate= 90.00,color=drawColor,anchor=base,inner sep=0pt, outer sep=0pt, scale=  0.90] at ( 43.20, 87.54) {0.2%
};

\node[rotate= 90.00,color=drawColor,anchor=base,inner sep=0pt, outer sep=0pt, scale=  0.90] at ( 43.20,119.50) {0.4%
};

\node[rotate= 90.00,color=drawColor,anchor=base,inner sep=0pt, outer sep=0pt, scale=  0.90] at ( 43.20,151.45) {0.6%
};

\node[rotate= 90.00,color=drawColor,anchor=base,inner sep=0pt, outer sep=0pt, scale=  0.90] at ( 43.20,183.40) {0.8%
};

\node[rotate= 90.00,color=drawColor,anchor=base,inner sep=0pt, outer sep=0pt, scale=  0.90] at ( 43.20,215.35) {1.0%
};

\draw[color=drawColor,line cap=round,line join=round,fill opacity=0.00,] ( 55.20, 49.20) --
    (250.54, 49.20) --
    (250.54,221.74) --
    ( 55.20,221.74) --
    ( 55.20, 49.20);
\end{scope}
\begin{scope}
\path[clip] (  0.00,  0.00) rectangle (252.94,252.94);
\definecolor[named]{drawColor}{rgb}{0.97,1.00,0.19}
\definecolor[named]{fillColor}{rgb}{0.88,0.08,0.52}
\definecolor[named]{drawColor}{rgb}{0.00,0.00,0.00}

\node[color=drawColor,anchor=base,inner sep=0pt, outer sep=0pt, scale=  1.00] at (152.87,  1.20) {$\gruen{\nz}$ resp. $\gruen{\nn}$%
};

\node[rotate= 90.00,color=drawColor,anchor=base,inner sep=0pt, outer sep=0pt, scale=  1.00] at ( 19.20,135.47) {$\rot{\yz}$ resp. $\rot{\yn}$%
};
\end{scope}
\begin{scope}
\path[clip] ( 55.20, 49.20) rectangle (250.54,221.74);
\definecolor[named]{drawColor}{rgb}{0.97,1.00,0.19}
\definecolor[named]{fillColor}{rgb}{0.88,0.08,0.52}
\definecolor[named]{drawColor}{rgb}{0.00,0.00,0.00}
\definecolor[named]{fillColor}{rgb}{0.75,0.75,0.75}

\draw[color=drawColor,line cap=round,line join=round,fill=fillColor,] ( 62.43,167.43) --
    ( 96.89,167.43) --
    ( 96.89,183.40) --
    ( 62.43,183.40) --
    cycle;
\end{scope}

}
\uncover<2>{
% Created by tikzDevice version 0.5.0 on 2011-07-19 12:37:03
\begin{scope}
\path[clip] ( 55.20, 49.20) rectangle (250.54,221.74);
\definecolor[named]{drawColor}{rgb}{0.00,0.13,0.19}
\definecolor[named]{fillColor}{rgb}{0.88,0.08,0.52}
\definecolor[named]{drawColor}{rgb}{0.75,0.75,0.75}
\definecolor[named]{fillColor}{rgb}{0.75,0.75,0.75}

\draw[color=drawColor,line cap=round,line join=round,fill=fillColor,] ( 96.89, 95.53) circle (  2.25);

\draw[color=drawColor,line cap=round,line join=round,fill=fillColor,] ( 96.89,175.41) circle (  2.25);

\draw[color=drawColor,line cap=round,line join=round,fill=fillColor,] (234.70,175.41) circle (  2.25);
\end{scope}
\begin{scope}
\path[clip] (  0.00,  0.00) rectangle (252.94,252.94);
\definecolor[named]{drawColor}{rgb}{0.00,0.13,0.19}
\definecolor[named]{fillColor}{rgb}{0.88,0.08,0.52}
\definecolor[named]{drawColor}{rgb}{0.00,0.00,0.00}

\draw[color=drawColor,line cap=round,line join=round,fill opacity=0.00,] ( 71.05, 49.20) -- (243.31, 49.20);

\draw[color=drawColor,line cap=round,line join=round,fill opacity=0.00,] ( 71.05, 49.20) -- ( 71.05, 43.20);

\draw[color=drawColor,line cap=round,line join=round,fill opacity=0.00,] (114.11, 49.20) -- (114.11, 43.20);

\draw[color=drawColor,line cap=round,line join=round,fill opacity=0.00,] (157.18, 49.20) -- (157.18, 43.20);

\draw[color=drawColor,line cap=round,line join=round,fill opacity=0.00,] (200.24, 49.20) -- (200.24, 43.20);

\draw[color=drawColor,line cap=round,line join=round,fill opacity=0.00,] (243.31, 49.20) -- (243.31, 43.20);

\node[color=drawColor,anchor=base,inner sep=0pt, outer sep=0pt, scale=  0.90] at ( 71.05, 25.20) {5%
};

\node[color=drawColor,anchor=base,inner sep=0pt, outer sep=0pt, scale=  0.90] at (114.11, 25.20) {10%
};

\node[color=drawColor,anchor=base,inner sep=0pt, outer sep=0pt, scale=  0.90] at (157.18, 25.20) {15%
};

\node[color=drawColor,anchor=base,inner sep=0pt, outer sep=0pt, scale=  0.90] at (200.24, 25.20) {20%
};

\node[color=drawColor,anchor=base,inner sep=0pt, outer sep=0pt, scale=  0.90] at (243.31, 25.20) {25%
};

\draw[color=drawColor,line cap=round,line join=round,fill opacity=0.00,] ( 55.20, 55.59) -- ( 55.20,215.35);

\draw[color=drawColor,line cap=round,line join=round,fill opacity=0.00,] ( 55.20, 55.59) -- ( 49.20, 55.59);

\draw[color=drawColor,line cap=round,line join=round,fill opacity=0.00,] ( 55.20, 87.54) -- ( 49.20, 87.54);

\draw[color=drawColor,line cap=round,line join=round,fill opacity=0.00,] ( 55.20,119.50) -- ( 49.20,119.50);

\draw[color=drawColor,line cap=round,line join=round,fill opacity=0.00,] ( 55.20,151.45) -- ( 49.20,151.45);

\draw[color=drawColor,line cap=round,line join=round,fill opacity=0.00,] ( 55.20,183.40) -- ( 49.20,183.40);

\draw[color=drawColor,line cap=round,line join=round,fill opacity=0.00,] ( 55.20,215.35) -- ( 49.20,215.35);

\node[rotate= 90.00,color=drawColor,anchor=base,inner sep=0pt, outer sep=0pt, scale=  0.90] at ( 43.20, 55.59) {0.0%
};

\node[rotate= 90.00,color=drawColor,anchor=base,inner sep=0pt, outer sep=0pt, scale=  0.90] at ( 43.20, 87.54) {0.2%
};

\node[rotate= 90.00,color=drawColor,anchor=base,inner sep=0pt, outer sep=0pt, scale=  0.90] at ( 43.20,119.50) {0.4%
};

\node[rotate= 90.00,color=drawColor,anchor=base,inner sep=0pt, outer sep=0pt, scale=  0.90] at ( 43.20,151.45) {0.6%
};

\node[rotate= 90.00,color=drawColor,anchor=base,inner sep=0pt, outer sep=0pt, scale=  0.90] at ( 43.20,183.40) {0.8%
};

\node[rotate= 90.00,color=drawColor,anchor=base,inner sep=0pt, outer sep=0pt, scale=  0.90] at ( 43.20,215.35) {1.0%
};

\draw[color=drawColor,line cap=round,line join=round,fill opacity=0.00,] ( 55.20, 49.20) --
    (250.54, 49.20) --
    (250.54,221.74) --
    ( 55.20,221.74) --
    ( 55.20, 49.20);
\end{scope}
\begin{scope}
\path[clip] (  0.00,  0.00) rectangle (252.94,252.94);
\definecolor[named]{drawColor}{rgb}{0.00,0.13,0.19}
\definecolor[named]{fillColor}{rgb}{0.88,0.08,0.52}
\definecolor[named]{drawColor}{rgb}{0.00,0.00,0.00}

\node[color=drawColor,anchor=base,inner sep=0pt, outer sep=0pt, scale=  1.00] at (152.87,  1.20) {$\gruen{\nz}$ resp. $\gruen{\nn}$%
};

\node[rotate= 90.00,color=drawColor,anchor=base,inner sep=0pt, outer sep=0pt, scale=  1.00] at ( 19.20,135.47) {$\rot{\yz}$ resp. $\rot{\yn}$%
};
\end{scope}
\begin{scope}
\path[clip] ( 55.20, 49.20) rectangle (250.54,221.74);
\definecolor[named]{drawColor}{rgb}{0.00,0.13,0.19}
\definecolor[named]{fillColor}{rgb}{0.88,0.08,0.52}
\definecolor[named]{drawColor}{rgb}{0.00,0.00,0.00}
\definecolor[named]{fillColor}{rgb}{0.75,0.75,0.75}

\draw[color=drawColor,line cap=round,line join=round,fill=fillColor,] ( 62.43,167.43) --
    ( 96.89,167.43) --
    ( 96.89,183.40) --
    ( 62.43,183.40) --
    cycle;

\draw[color=drawColor,line cap=round,line join=round,fill=fillColor,] (200.24,173.82) --
    (200.59,173.80) --
    (200.94,173.79) --
    (201.29,173.78) --
    (201.64,173.76) --
    (201.98,173.75) --
    (202.33,173.74) --
    (202.68,173.73) --
    (203.03,173.71) --
    (203.38,173.70) --
    (203.72,173.69) --
    (204.07,173.68) --
    (204.42,173.66) --
    (204.77,173.65) --
    (205.12,173.64) --
    (205.46,173.63) --
    (205.81,173.62) --
    (206.16,173.60) --
    (206.51,173.59) --
    (206.86,173.58) --
    (207.20,173.57) --
    (207.55,173.56) --
    (207.90,173.54) --
    (208.25,173.53) --
    (208.60,173.52) --
    (208.94,173.51) --
    (209.29,173.50) --
    (209.64,173.49) --
    (209.99,173.47) --
    (210.34,173.46) --
    (210.68,173.45) --
    (211.03,173.44) --
    (211.38,173.43) --
    (211.73,173.42) --
    (212.08,173.41) --
    (212.42,173.39) --
    (212.77,173.38) --
    (213.12,173.37) --
    (213.47,173.36) --
    (213.82,173.35) --
    (214.16,173.34) --
    (214.51,173.33) --
    (214.86,173.32) --
    (215.21,173.31) --
    (215.56,173.29) --
    (215.90,173.28) --
    (216.25,173.27) --
    (216.60,173.26) --
    (216.95,173.25) --
    (217.30,173.24) --
    (217.64,173.23) --
    (217.99,173.22) --
    (218.34,173.21) --
    (218.69,173.20) --
    (219.04,173.19) --
    (219.38,173.18) --
    (219.73,173.17) --
    (220.08,173.16) --
    (220.43,173.15) --
    (220.78,173.14) --
    (221.12,173.12) --
    (221.47,173.11) --
    (221.82,173.10) --
    (222.17,173.09) --
    (222.52,173.08) --
    (222.86,173.07) --
    (223.21,173.06) --
    (223.56,173.05) --
    (223.91,173.04) --
    (224.26,173.03) --
    (224.60,173.02) --
    (224.95,173.01) --
    (225.30,173.00) --
    (225.65,172.99) --
    (226.00,172.98) --
    (226.34,172.97) --
    (226.69,172.97) --
    (227.04,172.96) --
    (227.39,172.95) --
    (227.74,172.94) --
    (228.08,172.93) --
    (228.43,172.92) --
    (228.78,172.91) --
    (229.13,172.90) --
    (229.48,172.89) --
    (229.82,172.88) --
    (230.17,172.87) --
    (230.52,172.86) --
    (230.87,172.85) --
    (231.22,172.84) --
    (231.56,172.83) --
    (231.91,172.82) --
    (232.26,172.81) --
    (232.61,172.81) --
    (232.96,172.80) --
    (233.30,172.79) --
    (233.65,172.78) --
    (234.00,172.77) --
    (234.35,172.76) --
    (234.70,172.75) --
    (234.70,178.08) --
    (234.35,178.07) --
    (234.00,178.06) --
    (233.65,178.05) --
    (233.30,178.04) --
    (232.96,178.03) --
    (232.61,178.02) --
    (232.26,178.01) --
    (231.91,178.00) --
    (231.56,177.99) --
    (231.22,177.99) --
    (230.87,177.98) --
    (230.52,177.97) --
    (230.17,177.96) --
    (229.82,177.95) --
    (229.48,177.94) --
    (229.13,177.93) --
    (228.78,177.92) --
    (228.43,177.91) --
    (228.08,177.90) --
    (227.74,177.89) --
    (227.39,177.88) --
    (227.04,177.87) --
    (226.69,177.86) --
    (226.34,177.85) --
    (226.00,177.84) --
    (225.65,177.83) --
    (225.30,177.82) --
    (224.95,177.81) --
    (224.60,177.80) --
    (224.26,177.79) --
    (223.91,177.78) --
    (223.56,177.77) --
    (223.21,177.76) --
    (222.86,177.75) --
    (222.52,177.74) --
    (222.17,177.73) --
    (221.82,177.72) --
    (221.47,177.71) --
    (221.12,177.70) --
    (220.78,177.69) --
    (220.43,177.68) --
    (220.08,177.67) --
    (219.73,177.66) --
    (219.38,177.65) --
    (219.04,177.64) --
    (218.69,177.63) --
    (218.34,177.62) --
    (217.99,177.61) --
    (217.64,177.60) --
    (217.30,177.59) --
    (216.95,177.58) --
    (216.60,177.57) --
    (216.25,177.55) --
    (215.90,177.54) --
    (215.56,177.53) --
    (215.21,177.52) --
    (214.86,177.51) --
    (214.51,177.50) --
    (214.16,177.49) --
    (213.82,177.48) --
    (213.47,177.47) --
    (213.12,177.46) --
    (212.77,177.44) --
    (212.42,177.43) --
    (212.08,177.42) --
    (211.73,177.41) --
    (211.38,177.40) --
    (211.03,177.39) --
    (210.68,177.38) --
    (210.34,177.36) --
    (209.99,177.35) --
    (209.64,177.34) --
    (209.29,177.33) --
    (208.94,177.32) --
    (208.60,177.31) --
    (208.25,177.29) --
    (207.90,177.28) --
    (207.55,177.27) --
    (207.20,177.26) --
    (206.86,177.25) --
    (206.51,177.24) --
    (206.16,177.22) --
    (205.81,177.21) --
    (205.46,177.20) --
    (205.12,177.19) --
    (204.77,177.17) --
    (204.42,177.16) --
    (204.07,177.15) --
    (203.72,177.14) --
    (203.38,177.13) --
    (203.03,177.11) --
    (202.68,177.10) --
    (202.33,177.09) --
    (201.98,177.08) --
    (201.64,177.06) --
    (201.29,177.05) --
    (200.94,177.04) --
    (200.59,177.02) --
    (200.24,177.01) --
    cycle;

\draw[color=drawColor,line cap=round,line join=round,fill opacity=0.00,] (234.70,172.75) circle (  2.25);

\draw[color=drawColor,line cap=round,line join=round,fill opacity=0.00,] (234.70,178.08) circle (  2.25);
\end{scope}

\draw[-stealth,very thick] (110,175) -- (193,175) node [above,midway] {12 out of 16};
}
\uncover<3>{
% Created by tikzDevice version 0.5.0 on 2011-07-18 17:01:36
\begin{scope}
\path[clip] ( 55.20, 49.20) rectangle (250.54,221.74);
\definecolor[named]{drawColor}{rgb}{0.88,0.87,0.52}
\definecolor[named]{fillColor}{rgb}{0.88,0.08,0.52}
\definecolor[named]{drawColor}{rgb}{0.75,0.75,0.75}
\definecolor[named]{fillColor}{rgb}{0.75,0.75,0.75}

\draw[color=drawColor,line cap=round,line join=round,fill=fillColor,] ( 96.89, 95.53) circle (  2.25);

\draw[color=drawColor,line cap=round,line join=round,fill=fillColor,] ( 96.89,175.41) circle (  2.25);

\draw[color=drawColor,line cap=round,line join=round,fill=fillColor,] (234.70,175.41) circle (  2.25);
\end{scope}
\begin{scope}
\path[clip] (  0.00,  0.00) rectangle (252.94,252.94);
\definecolor[named]{drawColor}{rgb}{0.88,0.87,0.52}
\definecolor[named]{fillColor}{rgb}{0.88,0.08,0.52}
\definecolor[named]{drawColor}{rgb}{0.00,0.00,0.00}

\draw[color=drawColor,line cap=round,line join=round,fill opacity=0.00,] ( 71.05, 49.20) -- (243.31, 49.20);

\draw[color=drawColor,line cap=round,line join=round,fill opacity=0.00,] ( 71.05, 49.20) -- ( 71.05, 43.20);

\draw[color=drawColor,line cap=round,line join=round,fill opacity=0.00,] (114.11, 49.20) -- (114.11, 43.20);

\draw[color=drawColor,line cap=round,line join=round,fill opacity=0.00,] (157.18, 49.20) -- (157.18, 43.20);

\draw[color=drawColor,line cap=round,line join=round,fill opacity=0.00,] (200.24, 49.20) -- (200.24, 43.20);

\draw[color=drawColor,line cap=round,line join=round,fill opacity=0.00,] (243.31, 49.20) -- (243.31, 43.20);

\node[color=drawColor,anchor=base,inner sep=0pt, outer sep=0pt, scale=  0.90] at ( 71.05, 25.20) {5%
};

\node[color=drawColor,anchor=base,inner sep=0pt, outer sep=0pt, scale=  0.90] at (114.11, 25.20) {10%
};

\node[color=drawColor,anchor=base,inner sep=0pt, outer sep=0pt, scale=  0.90] at (157.18, 25.20) {15%
};

\node[color=drawColor,anchor=base,inner sep=0pt, outer sep=0pt, scale=  0.90] at (200.24, 25.20) {20%
};

\node[color=drawColor,anchor=base,inner sep=0pt, outer sep=0pt, scale=  0.90] at (243.31, 25.20) {25%
};

\draw[color=drawColor,line cap=round,line join=round,fill opacity=0.00,] ( 55.20, 55.59) -- ( 55.20,215.35);

\draw[color=drawColor,line cap=round,line join=round,fill opacity=0.00,] ( 55.20, 55.59) -- ( 49.20, 55.59);

\draw[color=drawColor,line cap=round,line join=round,fill opacity=0.00,] ( 55.20, 87.54) -- ( 49.20, 87.54);

\draw[color=drawColor,line cap=round,line join=round,fill opacity=0.00,] ( 55.20,119.50) -- ( 49.20,119.50);

\draw[color=drawColor,line cap=round,line join=round,fill opacity=0.00,] ( 55.20,151.45) -- ( 49.20,151.45);

\draw[color=drawColor,line cap=round,line join=round,fill opacity=0.00,] ( 55.20,183.40) -- ( 49.20,183.40);

\draw[color=drawColor,line cap=round,line join=round,fill opacity=0.00,] ( 55.20,215.35) -- ( 49.20,215.35);

\node[rotate= 90.00,color=drawColor,anchor=base,inner sep=0pt, outer sep=0pt, scale=  0.90] at ( 43.20, 55.59) {0.0%
};

\node[rotate= 90.00,color=drawColor,anchor=base,inner sep=0pt, outer sep=0pt, scale=  0.90] at ( 43.20, 87.54) {0.2%
};

\node[rotate= 90.00,color=drawColor,anchor=base,inner sep=0pt, outer sep=0pt, scale=  0.90] at ( 43.20,119.50) {0.4%
};

\node[rotate= 90.00,color=drawColor,anchor=base,inner sep=0pt, outer sep=0pt, scale=  0.90] at ( 43.20,151.45) {0.6%
};

\node[rotate= 90.00,color=drawColor,anchor=base,inner sep=0pt, outer sep=0pt, scale=  0.90] at ( 43.20,183.40) {0.8%
};

\node[rotate= 90.00,color=drawColor,anchor=base,inner sep=0pt, outer sep=0pt, scale=  0.90] at ( 43.20,215.35) {1.0%
};

\draw[color=drawColor,line cap=round,line join=round,fill opacity=0.00,] ( 55.20, 49.20) --
    (250.54, 49.20) --
    (250.54,221.74) --
    ( 55.20,221.74) --
    ( 55.20, 49.20);
\end{scope}
\begin{scope}
\path[clip] (  0.00,  0.00) rectangle (252.94,252.94);
\definecolor[named]{drawColor}{rgb}{0.88,0.87,0.52}
\definecolor[named]{fillColor}{rgb}{0.88,0.08,0.52}
\definecolor[named]{drawColor}{rgb}{0.00,0.00,0.00}

\node[color=drawColor,anchor=base,inner sep=0pt, outer sep=0pt, scale=  1.00] at (152.87,  1.20) {$\gruen{\nz}$ resp. $\gruen{\nn}$%
};

\node[rotate= 90.00,color=drawColor,anchor=base,inner sep=0pt, outer sep=0pt, scale=  1.00] at ( 19.20,135.47) {$\rot{\yz}$ resp. $\rot{\yn}$%
};
\end{scope}
\begin{scope}
\path[clip] ( 55.20, 49.20) rectangle (250.54,221.74);
\definecolor[named]{drawColor}{rgb}{0.88,0.87,0.52}
\definecolor[named]{fillColor}{rgb}{0.88,0.08,0.52}
\definecolor[named]{drawColor}{rgb}{0.00,0.00,0.00}
\definecolor[named]{fillColor}{rgb}{0.75,0.75,0.75}

\draw[color=drawColor,line cap=round,line join=round,fill=fillColor,] ( 62.43,167.43) --
    ( 96.89,167.43) --
    ( 96.89,183.40) --
    ( 62.43,183.40) --
    cycle;

\draw[color=drawColor,line cap=round,line join=round,fill=fillColor,] (200.24,173.82) --
    (200.59,173.80) --
    (200.94,173.79) --
    (201.29,173.78) --
    (201.64,173.76) --
    (201.98,173.75) --
    (202.33,173.74) --
    (202.68,173.73) --
    (203.03,173.71) --
    (203.38,173.70) --
    (203.72,173.69) --
    (204.07,173.68) --
    (204.42,173.66) --
    (204.77,173.65) --
    (205.12,173.64) --
    (205.46,173.63) --
    (205.81,173.62) --
    (206.16,173.60) --
    (206.51,173.59) --
    (206.86,173.58) --
    (207.20,173.57) --
    (207.55,173.56) --
    (207.90,173.54) --
    (208.25,173.53) --
    (208.60,173.52) --
    (208.94,173.51) --
    (209.29,173.50) --
    (209.64,173.49) --
    (209.99,173.47) --
    (210.34,173.46) --
    (210.68,173.45) --
    (211.03,173.44) --
    (211.38,173.43) --
    (211.73,173.42) --
    (212.08,173.41) --
    (212.42,173.39) --
    (212.77,173.38) --
    (213.12,173.37) --
    (213.47,173.36) --
    (213.82,173.35) --
    (214.16,173.34) --
    (214.51,173.33) --
    (214.86,173.32) --
    (215.21,173.31) --
    (215.56,173.29) --
    (215.90,173.28) --
    (216.25,173.27) --
    (216.60,173.26) --
    (216.95,173.25) --
    (217.30,173.24) --
    (217.64,173.23) --
    (217.99,173.22) --
    (218.34,173.21) --
    (218.69,173.20) --
    (219.04,173.19) --
    (219.38,173.18) --
    (219.73,173.17) --
    (220.08,173.16) --
    (220.43,173.15) --
    (220.78,173.14) --
    (221.12,173.12) --
    (221.47,173.11) --
    (221.82,173.10) --
    (222.17,173.09) --
    (222.52,173.08) --
    (222.86,173.07) --
    (223.21,173.06) --
    (223.56,173.05) --
    (223.91,173.04) --
    (224.26,173.03) --
    (224.60,173.02) --
    (224.95,173.01) --
    (225.30,173.00) --
    (225.65,172.99) --
    (226.00,172.98) --
    (226.34,172.97) --
    (226.69,172.97) --
    (227.04,172.96) --
    (227.39,172.95) --
    (227.74,172.94) --
    (228.08,172.93) --
    (228.43,172.92) --
    (228.78,172.91) --
    (229.13,172.90) --
    (229.48,172.89) --
    (229.82,172.88) --
    (230.17,172.87) --
    (230.52,172.86) --
    (230.87,172.85) --
    (231.22,172.84) --
    (231.56,172.83) --
    (231.91,172.82) --
    (232.26,172.81) --
    (232.61,172.81) --
    (232.96,172.80) --
    (233.30,172.79) --
    (233.65,172.78) --
    (234.00,172.77) --
    (234.35,172.76) --
    (234.70,172.75) --
    (234.70,178.08) --
    (234.35,178.07) --
    (234.00,178.06) --
    (233.65,178.05) --
    (233.30,178.04) --
    (232.96,178.03) --
    (232.61,178.02) --
    (232.26,178.01) --
    (231.91,178.00) --
    (231.56,177.99) --
    (231.22,177.99) --
    (230.87,177.98) --
    (230.52,177.97) --
    (230.17,177.96) --
    (229.82,177.95) --
    (229.48,177.94) --
    (229.13,177.93) --
    (228.78,177.92) --
    (228.43,177.91) --
    (228.08,177.90) --
    (227.74,177.89) --
    (227.39,177.88) --
    (227.04,177.87) --
    (226.69,177.86) --
    (226.34,177.85) --
    (226.00,177.84) --
    (225.65,177.83) --
    (225.30,177.82) --
    (224.95,177.81) --
    (224.60,177.80) --
    (224.26,177.79) --
    (223.91,177.78) --
    (223.56,177.77) --
    (223.21,177.76) --
    (222.86,177.75) --
    (222.52,177.74) --
    (222.17,177.73) --
    (221.82,177.72) --
    (221.47,177.71) --
    (221.12,177.70) --
    (220.78,177.69) --
    (220.43,177.68) --
    (220.08,177.67) --
    (219.73,177.66) --
    (219.38,177.65) --
    (219.04,177.64) --
    (218.69,177.63) --
    (218.34,177.62) --
    (217.99,177.61) --
    (217.64,177.60) --
    (217.30,177.59) --
    (216.95,177.58) --
    (216.60,177.57) --
    (216.25,177.55) --
    (215.90,177.54) --
    (215.56,177.53) --
    (215.21,177.52) --
    (214.86,177.51) --
    (214.51,177.50) --
    (214.16,177.49) --
    (213.82,177.48) --
    (213.47,177.47) --
    (213.12,177.46) --
    (212.77,177.44) --
    (212.42,177.43) --
    (212.08,177.42) --
    (211.73,177.41) --
    (211.38,177.40) --
    (211.03,177.39) --
    (210.68,177.38) --
    (210.34,177.36) --
    (209.99,177.35) --
    (209.64,177.34) --
    (209.29,177.33) --
    (208.94,177.32) --
    (208.60,177.31) --
    (208.25,177.29) --
    (207.90,177.28) --
    (207.55,177.27) --
    (207.20,177.26) --
    (206.86,177.25) --
    (206.51,177.24) --
    (206.16,177.22) --
    (205.81,177.21) --
    (205.46,177.20) --
    (205.12,177.19) --
    (204.77,177.17) --
    (204.42,177.16) --
    (204.07,177.15) --
    (203.72,177.14) --
    (203.38,177.13) --
    (203.03,177.11) --
    (202.68,177.10) --
    (202.33,177.09) --
    (201.98,177.08) --
    (201.64,177.06) --
    (201.29,177.05) --
    (200.94,177.04) --
    (200.59,177.02) --
    (200.24,177.01) --
    cycle;

\draw[color=drawColor,line cap=round,line join=round,fill=fillColor,] ( 62.43, 87.54) --
    ( 96.89, 87.54) --
    ( 96.89,103.52) --
    ( 62.43,103.52) --
    cycle;
\definecolor[named]{drawColor}{rgb}{0.75,0.75,0.75}

\draw[color=drawColor,line cap=round,line join=round,fill opacity=0.00,] (234.70,172.75) circle (  2.25);

\draw[color=drawColor,line cap=round,line join=round,fill opacity=0.00,] (234.70,178.08) circle (  2.25);
\end{scope}

\draw[-stealth,very thick,gray] (110,175) -- (193,175) node [above,midway,gray] {12 out of 16};
}
\uncover<4>{
% Created by tikzDevice version 0.5.0 on 2011-07-18 17:01:45
\begin{scope}
\path[clip] ( 55.20, 49.20) rectangle (250.54,221.74);
\definecolor[named]{drawColor}{rgb}{0.60,0.10,0.52}
\definecolor[named]{fillColor}{rgb}{0.88,0.08,0.52}
\definecolor[named]{drawColor}{rgb}{0.75,0.75,0.75}
\definecolor[named]{fillColor}{rgb}{0.75,0.75,0.75}

\draw[color=drawColor,line cap=round,line join=round,fill=fillColor,] ( 96.89, 95.53) circle (  2.25);

\draw[color=drawColor,line cap=round,line join=round,fill=fillColor,] ( 96.89,175.41) circle (  2.25);

\draw[color=drawColor,line cap=round,line join=round,fill=fillColor,] (234.70,175.41) circle (  2.25);
\end{scope}
\begin{scope}
\path[clip] (  0.00,  0.00) rectangle (252.94,252.94);
\definecolor[named]{drawColor}{rgb}{0.60,0.10,0.52}
\definecolor[named]{fillColor}{rgb}{0.88,0.08,0.52}
\definecolor[named]{drawColor}{rgb}{0.00,0.00,0.00}

\draw[color=drawColor,line cap=round,line join=round,fill opacity=0.00,] ( 71.05, 49.20) -- (243.31, 49.20);

\draw[color=drawColor,line cap=round,line join=round,fill opacity=0.00,] ( 71.05, 49.20) -- ( 71.05, 43.20);

\draw[color=drawColor,line cap=round,line join=round,fill opacity=0.00,] (114.11, 49.20) -- (114.11, 43.20);

\draw[color=drawColor,line cap=round,line join=round,fill opacity=0.00,] (157.18, 49.20) -- (157.18, 43.20);

\draw[color=drawColor,line cap=round,line join=round,fill opacity=0.00,] (200.24, 49.20) -- (200.24, 43.20);

\draw[color=drawColor,line cap=round,line join=round,fill opacity=0.00,] (243.31, 49.20) -- (243.31, 43.20);

\node[color=drawColor,anchor=base,inner sep=0pt, outer sep=0pt, scale=  0.90] at ( 71.05, 25.20) {5%
};

\node[color=drawColor,anchor=base,inner sep=0pt, outer sep=0pt, scale=  0.90] at (114.11, 25.20) {10%
};

\node[color=drawColor,anchor=base,inner sep=0pt, outer sep=0pt, scale=  0.90] at (157.18, 25.20) {15%
};

\node[color=drawColor,anchor=base,inner sep=0pt, outer sep=0pt, scale=  0.90] at (200.24, 25.20) {20%
};

\node[color=drawColor,anchor=base,inner sep=0pt, outer sep=0pt, scale=  0.90] at (243.31, 25.20) {25%
};

\draw[color=drawColor,line cap=round,line join=round,fill opacity=0.00,] ( 55.20, 55.59) -- ( 55.20,215.35);

\draw[color=drawColor,line cap=round,line join=round,fill opacity=0.00,] ( 55.20, 55.59) -- ( 49.20, 55.59);

\draw[color=drawColor,line cap=round,line join=round,fill opacity=0.00,] ( 55.20, 87.54) -- ( 49.20, 87.54);

\draw[color=drawColor,line cap=round,line join=round,fill opacity=0.00,] ( 55.20,119.50) -- ( 49.20,119.50);

\draw[color=drawColor,line cap=round,line join=round,fill opacity=0.00,] ( 55.20,151.45) -- ( 49.20,151.45);

\draw[color=drawColor,line cap=round,line join=round,fill opacity=0.00,] ( 55.20,183.40) -- ( 49.20,183.40);

\draw[color=drawColor,line cap=round,line join=round,fill opacity=0.00,] ( 55.20,215.35) -- ( 49.20,215.35);

\node[rotate= 90.00,color=drawColor,anchor=base,inner sep=0pt, outer sep=0pt, scale=  0.90] at ( 43.20, 55.59) {0.0%
};

\node[rotate= 90.00,color=drawColor,anchor=base,inner sep=0pt, outer sep=0pt, scale=  0.90] at ( 43.20, 87.54) {0.2%
};

\node[rotate= 90.00,color=drawColor,anchor=base,inner sep=0pt, outer sep=0pt, scale=  0.90] at ( 43.20,119.50) {0.4%
};

\node[rotate= 90.00,color=drawColor,anchor=base,inner sep=0pt, outer sep=0pt, scale=  0.90] at ( 43.20,151.45) {0.6%
};

\node[rotate= 90.00,color=drawColor,anchor=base,inner sep=0pt, outer sep=0pt, scale=  0.90] at ( 43.20,183.40) {0.8%
};

\node[rotate= 90.00,color=drawColor,anchor=base,inner sep=0pt, outer sep=0pt, scale=  0.90] at ( 43.20,215.35) {1.0%
};

\draw[color=drawColor,line cap=round,line join=round,fill opacity=0.00,] ( 55.20, 49.20) --
    (250.54, 49.20) --
    (250.54,221.74) --
    ( 55.20,221.74) --
    ( 55.20, 49.20);
\end{scope}
\begin{scope}
\path[clip] (  0.00,  0.00) rectangle (252.94,252.94);
\definecolor[named]{drawColor}{rgb}{0.60,0.10,0.52}
\definecolor[named]{fillColor}{rgb}{0.88,0.08,0.52}
\definecolor[named]{drawColor}{rgb}{0.00,0.00,0.00}

\node[color=drawColor,anchor=base,inner sep=0pt, outer sep=0pt, scale=  1.00] at (152.87,  1.20) {$\gruen{\nz}$ resp. $\gruen{\nn}$%
};

\node[rotate= 90.00,color=drawColor,anchor=base,inner sep=0pt, outer sep=0pt, scale=  1.00] at ( 19.20,135.47) {$\rot{\yz}$ resp. $\rot{\yn}$%
};
\end{scope}
\begin{scope}
\path[clip] ( 55.20, 49.20) rectangle (250.54,221.74);
\definecolor[named]{drawColor}{rgb}{0.60,0.10,0.52}
\definecolor[named]{fillColor}{rgb}{0.88,0.08,0.52}
\definecolor[named]{drawColor}{rgb}{0.00,0.00,0.00}
\definecolor[named]{fillColor}{rgb}{0.75,0.75,0.75}

\draw[color=drawColor,line cap=round,line join=round,fill=fillColor,] ( 62.43,167.43) --
    ( 96.89,167.43) --
    ( 96.89,183.40) --
    ( 62.43,183.40) --
    cycle;

\draw[color=drawColor,line cap=round,line join=round,fill=fillColor,] (200.24,173.82) --
    (200.59,173.80) --
    (200.94,173.79) --
    (201.29,173.78) --
    (201.64,173.76) --
    (201.98,173.75) --
    (202.33,173.74) --
    (202.68,173.73) --
    (203.03,173.71) --
    (203.38,173.70) --
    (203.72,173.69) --
    (204.07,173.68) --
    (204.42,173.66) --
    (204.77,173.65) --
    (205.12,173.64) --
    (205.46,173.63) --
    (205.81,173.62) --
    (206.16,173.60) --
    (206.51,173.59) --
    (206.86,173.58) --
    (207.20,173.57) --
    (207.55,173.56) --
    (207.90,173.54) --
    (208.25,173.53) --
    (208.60,173.52) --
    (208.94,173.51) --
    (209.29,173.50) --
    (209.64,173.49) --
    (209.99,173.47) --
    (210.34,173.46) --
    (210.68,173.45) --
    (211.03,173.44) --
    (211.38,173.43) --
    (211.73,173.42) --
    (212.08,173.41) --
    (212.42,173.39) --
    (212.77,173.38) --
    (213.12,173.37) --
    (213.47,173.36) --
    (213.82,173.35) --
    (214.16,173.34) --
    (214.51,173.33) --
    (214.86,173.32) --
    (215.21,173.31) --
    (215.56,173.29) --
    (215.90,173.28) --
    (216.25,173.27) --
    (216.60,173.26) --
    (216.95,173.25) --
    (217.30,173.24) --
    (217.64,173.23) --
    (217.99,173.22) --
    (218.34,173.21) --
    (218.69,173.20) --
    (219.04,173.19) --
    (219.38,173.18) --
    (219.73,173.17) --
    (220.08,173.16) --
    (220.43,173.15) --
    (220.78,173.14) --
    (221.12,173.12) --
    (221.47,173.11) --
    (221.82,173.10) --
    (222.17,173.09) --
    (222.52,173.08) --
    (222.86,173.07) --
    (223.21,173.06) --
    (223.56,173.05) --
    (223.91,173.04) --
    (224.26,173.03) --
    (224.60,173.02) --
    (224.95,173.01) --
    (225.30,173.00) --
    (225.65,172.99) --
    (226.00,172.98) --
    (226.34,172.97) --
    (226.69,172.97) --
    (227.04,172.96) --
    (227.39,172.95) --
    (227.74,172.94) --
    (228.08,172.93) --
    (228.43,172.92) --
    (228.78,172.91) --
    (229.13,172.90) --
    (229.48,172.89) --
    (229.82,172.88) --
    (230.17,172.87) --
    (230.52,172.86) --
    (230.87,172.85) --
    (231.22,172.84) --
    (231.56,172.83) --
    (231.91,172.82) --
    (232.26,172.81) --
    (232.61,172.81) --
    (232.96,172.80) --
    (233.30,172.79) --
    (233.65,172.78) --
    (234.00,172.77) --
    (234.35,172.76) --
    (234.70,172.75) --
    (234.70,178.08) --
    (234.35,178.07) --
    (234.00,178.06) --
    (233.65,178.05) --
    (233.30,178.04) --
    (232.96,178.03) --
    (232.61,178.02) --
    (232.26,178.01) --
    (231.91,178.00) --
    (231.56,177.99) --
    (231.22,177.99) --
    (230.87,177.98) --
    (230.52,177.97) --
    (230.17,177.96) --
    (229.82,177.95) --
    (229.48,177.94) --
    (229.13,177.93) --
    (228.78,177.92) --
    (228.43,177.91) --
    (228.08,177.90) --
    (227.74,177.89) --
    (227.39,177.88) --
    (227.04,177.87) --
    (226.69,177.86) --
    (226.34,177.85) --
    (226.00,177.84) --
    (225.65,177.83) --
    (225.30,177.82) --
    (224.95,177.81) --
    (224.60,177.80) --
    (224.26,177.79) --
    (223.91,177.78) --
    (223.56,177.77) --
    (223.21,177.76) --
    (222.86,177.75) --
    (222.52,177.74) --
    (222.17,177.73) --
    (221.82,177.72) --
    (221.47,177.71) --
    (221.12,177.70) --
    (220.78,177.69) --
    (220.43,177.68) --
    (220.08,177.67) --
    (219.73,177.66) --
    (219.38,177.65) --
    (219.04,177.64) --
    (218.69,177.63) --
    (218.34,177.62) --
    (217.99,177.61) --
    (217.64,177.60) --
    (217.30,177.59) --
    (216.95,177.58) --
    (216.60,177.57) --
    (216.25,177.55) --
    (215.90,177.54) --
    (215.56,177.53) --
    (215.21,177.52) --
    (214.86,177.51) --
    (214.51,177.50) --
    (214.16,177.49) --
    (213.82,177.48) --
    (213.47,177.47) --
    (213.12,177.46) --
    (212.77,177.44) --
    (212.42,177.43) --
    (212.08,177.42) --
    (211.73,177.41) --
    (211.38,177.40) --
    (211.03,177.39) --
    (210.68,177.38) --
    (210.34,177.36) --
    (209.99,177.35) --
    (209.64,177.34) --
    (209.29,177.33) --
    (208.94,177.32) --
    (208.60,177.31) --
    (208.25,177.29) --
    (207.90,177.28) --
    (207.55,177.27) --
    (207.20,177.26) --
    (206.86,177.25) --
    (206.51,177.24) --
    (206.16,177.22) --
    (205.81,177.21) --
    (205.46,177.20) --
    (205.12,177.19) --
    (204.77,177.17) --
    (204.42,177.16) --
    (204.07,177.15) --
    (203.72,177.14) --
    (203.38,177.13) --
    (203.03,177.11) --
    (202.68,177.10) --
    (202.33,177.09) --
    (201.98,177.08) --
    (201.64,177.06) --
    (201.29,177.05) --
    (200.94,177.04) --
    (200.59,177.02) --
    (200.24,177.01) --
    cycle;

\draw[color=drawColor,line cap=round,line join=round,fill=fillColor,] ( 62.43, 87.54) --
    ( 96.89, 87.54) --
    ( 96.89,103.52) --
    ( 62.43,103.52) --
    cycle;
\definecolor[named]{drawColor}{rgb}{0.75,0.75,0.75}

\draw[color=drawColor,line cap=round,line join=round,fill opacity=0.00,] (234.70,172.75) circle (  2.25);

\draw[color=drawColor,line cap=round,line join=round,fill opacity=0.00,] (234.70,178.08) circle (  2.25);
\definecolor[named]{drawColor}{rgb}{0.00,0.00,0.00}

\draw[color=drawColor,line cap=round,line join=round,fill=fillColor,] (200.24,189.79) --
    (200.59,189.59) --
    (200.94,189.38) --
    (201.29,189.18) --
    (201.64,188.97) --
    (201.98,188.77) --
    (202.33,188.57) --
    (202.68,188.37) --
    (203.03,188.17) --
    (203.38,187.97) --
    (203.72,187.77) --
    (204.07,187.57) --
    (204.42,187.37) --
    (204.77,187.18) --
    (205.12,186.98) --
    (205.46,186.78) --
    (205.81,186.59) --
    (206.16,186.40) --
    (206.51,186.20) --
    (206.86,186.01) --
    (207.20,185.82) --
    (207.55,185.63) --
    (207.90,185.44) --
    (208.25,185.25) --
    (208.60,185.06) --
    (208.94,184.88) --
    (209.29,184.69) --
    (209.64,184.50) --
    (209.99,184.32) --
    (210.34,184.13) --
    (210.68,183.95) --
    (211.03,183.77) --
    (211.38,183.58) --
    (211.73,183.40) --
    (212.08,183.22) --
    (212.42,183.04) --
    (212.77,182.86) --
    (213.12,182.68) --
    (213.47,182.50) --
    (213.82,182.32) --
    (214.16,182.15) --
    (214.51,181.97) --
    (214.86,181.80) --
    (215.21,181.62) --
    (215.56,181.45) --
    (215.90,181.27) --
    (216.25,181.10) --
    (216.60,180.93) --
    (216.95,180.75) --
    (217.30,180.58) --
    (217.64,180.41) --
    (217.99,180.24) --
    (218.34,180.07) --
    (218.69,179.90) --
    (219.04,179.73) --
    (219.38,179.57) --
    (219.73,179.40) --
    (220.08,179.23) --
    (220.43,179.07) --
    (220.78,178.90) --
    (221.12,178.74) --
    (221.47,178.57) --
    (221.82,178.41) --
    (222.17,178.25) --
    (222.52,178.09) --
    (222.86,177.92) --
    (223.21,177.76) --
    (223.56,177.60) --
    (223.91,177.44) --
    (224.26,177.28) --
    (224.60,177.12) --
    (224.95,176.97) --
    (225.30,176.81) --
    (225.65,176.65) --
    (226.00,176.49) --
    (226.34,176.34) --
    (226.69,176.18) --
    (227.04,176.03) --
    (227.39,175.87) --
    (227.74,175.72) --
    (228.08,175.57) --
    (228.43,175.41) --
    (228.78,175.26) --
    (229.13,175.11) --
    (229.48,174.96) --
    (229.82,174.81) --
    (230.17,174.66) --
    (230.52,174.51) --
    (230.87,174.36) --
    (231.22,174.21) --
    (231.56,174.06) --
    (231.91,173.91) --
    (232.26,173.77) --
    (232.61,173.62) --
    (232.96,173.47) --
    (233.30,173.33) --
    (233.65,173.18) --
    (234.00,173.04) --
    (234.35,172.89) --
    (234.70,172.75) --
    (234.70,178.08) --
    (234.35,178.20) --
    (234.00,178.33) --
    (233.65,178.45) --
    (233.30,178.58) --
    (232.96,178.71) --
    (232.61,178.84) --
    (232.26,178.97) --
    (231.91,179.09) --
    (231.56,179.22) --
    (231.22,179.35) --
    (230.87,179.48) --
    (230.52,179.61) --
    (230.17,179.74) --
    (229.82,179.88) --
    (229.48,180.01) --
    (229.13,180.14) --
    (228.78,180.27) --
    (228.43,180.41) --
    (228.08,180.54) --
    (227.74,180.67) --
    (227.39,180.81) --
    (227.04,180.94) --
    (226.69,181.08) --
    (226.34,181.22) --
    (226.00,181.35) --
    (225.65,181.49) --
    (225.30,181.63) --
    (224.95,181.76) --
    (224.60,181.90) --
    (224.26,182.04) --
    (223.91,182.18) --
    (223.56,182.32) --
    (223.21,182.46) --
    (222.86,182.60) --
    (222.52,182.74) --
    (222.17,182.89) --
    (221.82,183.03) --
    (221.47,183.17) --
    (221.12,183.32) --
    (220.78,183.46) --
    (220.43,183.60) --
    (220.08,183.75) --
    (219.73,183.89) --
    (219.38,184.04) --
    (219.04,184.19) --
    (218.69,184.33) --
    (218.34,184.48) --
    (217.99,184.63) --
    (217.64,184.78) --
    (217.30,184.93) --
    (216.95,185.08) --
    (216.60,185.23) --
    (216.25,185.38) --
    (215.90,185.53) --
    (215.56,185.68) --
    (215.21,185.84) --
    (214.86,185.99) --
    (214.51,186.14) --
    (214.16,186.30) --
    (213.82,186.45) --
    (213.47,186.61) --
    (213.12,186.77) --
    (212.77,186.92) --
    (212.42,187.08) --
    (212.08,187.24) --
    (211.73,187.40) --
    (211.38,187.55) --
    (211.03,187.71) --
    (210.68,187.88) --
    (210.34,188.04) --
    (209.99,188.20) --
    (209.64,188.36) --
    (209.29,188.52) --
    (208.94,188.69) --
    (208.60,188.85) --
    (208.25,189.01) --
    (207.90,189.18) --
    (207.55,189.35) --
    (207.20,189.51) --
    (206.86,189.68) --
    (206.51,189.85) --
    (206.16,190.02) --
    (205.81,190.19) --
    (205.46,190.36) --
    (205.12,190.53) --
    (204.77,190.70) --
    (204.42,190.87) --
    (204.07,191.04) --
    (203.72,191.22) --
    (203.38,191.39) --
    (203.03,191.56) --
    (202.68,191.74) --
    (202.33,191.92) --
    (201.98,192.09) --
    (201.64,192.27) --
    (201.29,192.45) --
    (200.94,192.63) --
    (200.59,192.81) --
    (200.24,192.99) --
    cycle;

\draw[color=drawColor,line cap=round,line join=round,fill opacity=0.00,] (234.70,172.75) circle (  2.25);

\draw[color=drawColor,line cap=round,line join=round,fill opacity=0.00,] (200.24,192.99) circle (  2.25);
\end{scope}

\draw[-stealth,very thick,gray] (110,175) -- (193,175) node [above,midway,gray] {12 out of 16};
\draw[-stealth,very thick] (110,102) -- (213,180) node [above,midway,sloped] {16 out of 16};
%\draw[-stealth,very thick, postaction={decorate,decoration={raise=2pt,text along path, %
%                                       text={16 out of 16}}}] (110,95) .. controls +(right:30) and +(down:30) .. (210,183);
%\draw[-stealth,very thick] (110,95) .. controls +(right:10) and +(down:10) .. (210,183) node [below,midway,sloped] {16 out of 16};
}
\end{tikzpicture}
\end{column}
\begin{column}{0.45\textwidth}
\uncover<1->{%
\begin{block}{no conflict:}
prior $\nzg\! \in [4, 8]$, $\yzr\! \in [0.7, 0.8]$\\
data $s/n = 12/16 = 0.75$
\end{block}
} %
\uncover<2->{%
\vspace*{-1.5ex}\centerline{\color{lmugreen} $\blacktriangledown$}\vspace*{-1.5ex}
\begin{block}{}
%$\gruen{\nn} \in [20, 24]$,
\centering $\ynr \in [0.73, 0.77]$
%\centering ``spotlight'' shape
\end{block}
} %
\uncover<3->{%
\begin{block}{\pdc:}
prior $\nzg \in\! [4, 8]$, $\yzr\! \in [0.2, 0.3]$\\
data $s/n = 16/16 = 1$
\end{block}
} %
\uncover<4->{%
\vspace*{-1.5ex}\centerline{\color{lmugreen} $\blacktriangledown$}\vspace*{-1.5ex}
\begin{block}{}
%$\gruen{\nn} \in [20, 24]$,
\centering $\ynr \in [0.73, 0.86]$
%\centering ``banana'' shape
\end{block}
} %
\end{column}
\end{columns}

}

\begin{frame}{Scaled Normal Data Example ($\x \stackrel{iid}{\sim} \norm(\mu,1)$)}
\includegraphics[trim = 20mm 20mm 20mm 25mm, clip, width=\textwidth]{../fig/jstp-paper_nv_nvar_vertfu01_081128}
\end{frame}

\begin{frame}[label=othermodels-back]{General Model Properties}

\uncover<1->{%
Favourable inference properties \hyperlink{othermodels-app}{\small (cf.\ other models based on sets of priors)}}
\begin{itemize}%[<+->]
\note<10>[item]{In contrast to other models based on sets of priors,\\
this model is easy to handle and to elicit,\\ and has very favourable inference properties}
\item<1-> $n \to \infty$\hspace*{0.9ex}%
%\uncover<2->{\then\ \ $\ynr$ values in $\PNc \to \frac{\tau(\x)}{n}$\hspace*{0.9ex}}%
%\uncover<3->{\then\ \ `Bayesian consistency'}
\uncover<2->{\then\ \ $\ynr$ stretch in $\PNc \to 0$\hspace*{0.9ex}}%
\uncover<3->{\then\ \ precise inferences}
\note<10>[item]{`Bayesian consistency' just like in usual Bayesian inference}
\item<4-> larger $\nzg$ \then\ larger $\PNc$ \then\ more vague inferences
\note<10>[item]{larger $\nzg$ as compared to $n$:
more weight on imprecise prior $\MZ$\\ leads to more imprecise posterior $\MN$}
\item<5-> larger range of $\yzr$ in $\PZc$ \then\ larger range of $\ynr$ in $\PNc$\\ \quad  \then\ more vague inferences
\note<10>[item]{more imprecise prior $\MZ$ leads to more imprecise $\MN$}
\end{itemize}

\uncover<6->{%
Model very easy to handle:}
\begin{itemize}%[<+->]
\item<6-> Hyperparameter set $\PZc$ defines set of priors $\MZ$
\item<7-> Hyperparameter set $\PNc$ defines set of posteriors $\MN$
\item<8-> $\PZc \to \PNc$ is easy: $\nng = \nzg + n$, $\ynr = \frac{\nzg}{\nzg + n}\yzr + \frac{n}{\nzg + n}\frac{\tau(\x)}{n}$
\item<9-> For quantities linear in posteriors, bounds are attained at ``pure'' posteriors $p(\psib\mid\nng,\bynr)$
\then\ optimise over $\PNc$ (not over $\MN$)
\item<10-> Often, optimising over $(\nng,\bynr) \in \PNc$ is also easy:\\
\alert{closed form solution} for $\ynr$ = posterior `guess' for $\frac{\tau(\x)}{n}$ (think: $\bar{\x}$)\\
given $\PZc$ has `nice' shape
\end{itemize}

\end{frame}


\subsection{Parameter Set Shapes}

\frame{\frametitle{Parameter Set Shapes}

%graphs line, rectangle, eggplant
\vspace*{-7ex}
\begin{tikzpicture}
%\uncover<1>{\node {\includegraphics[scale=0.75]{../R/shape0.pdf}};
%            \draw[-stealth,very thick] (-1,-0.4) -- (2.3,1.1) node [above,midway,sloped] {+ data};}
\uncover<1>{\node {\includegraphics[scale=0.75]{../R/shape1.pdf}};
            \draw[-stealth,very thick] (-1,-0.4) -- (2.3,1.1) node [above,midway,sloped] {+ data};}
\uncover<2>{\node {\includegraphics[scale=0.75]{../R/shape2.pdf}};
            \draw[-stealth,very thick] ( 0,-0.5) -- (2.3,1.1) node [above,midway,sloped] {+ data};}
\uncover<3>{\node {\includegraphics[scale=0.75]{../R/shape3.pdf}};
            \draw[-stealth,very thick] ( 0,-0.5) -- (2.3,1.1) node [above,midway,sloped] {+ data};}
\end{tikzpicture}
}


\begin{frame}[label=luck-back]{Parameter Set Shapes}
\begin{itemize}%[<+->]
\item<1-> Set shape is crucial modeling choice:\\ %A range of $\nzg$ values is needed for \pdc\ reaction
trade-off between model complexity and model behaviour
% \begin{itemize}
\item<2-> $\PZc\! = \nzg \times [\yzlr, \yzur]$ {\small (Walley\,1996;\,Quaghebeur\,\&\,de\,Cooman\,2005)}:\\
 $\PNc\! = \nng \times [\ynlr, \ynur]$ \then\ optimise over $[\ynlr, \ynur]$ only,\\
 but no \pdc\ sensitivity
\item<3-> $\PZc\! = [\nzlg,\nzug] \times [\yzlr, \yzur]$ {\small (Walley 1991; Walter \& Augustin 2009)}:\\
 $\PNc$ have non-trivial forms (banana / spotlight), but \pdc\ sensitivity and closed form for $\min / \max \ynr$ over $\PNc$.\\
 For other inferences, \hyperlink{luck-app}{\textbf{R} package \texttt{luck}} implements optimisation
 over $\PNc$ via box-constraint optimisation over $\PZc$
% \end{itemize}
\item<4-> Other set shapes are possible, but may be more difficult to handle %elicit
%\item<8-> Prior information may be such that range of $\yzr$ changes with $\nzg$ (or vice versa)
\end{itemize}
\end{frame}

\begin{frame}[label=boat-back]{Parameter Set Shape for Strong Prior-Data Agreement}
%\hspace*{-2ex}%
%\uncover<2->{%
%\vspace*{-7ex}%
%\begin{center} %trim=l b r t
\hyperlink{boat-app}{%
\begin{tikzpicture}
%\includegraphics[trim = 40mm 35mm 40mm 45mm, clip, width=\textwidth]{../R/boatshape-durham0913} % width=8, height=6
\uncover<1>{\node {\includegraphics[trim = 70mm 52.5mm 70mm 65mm, clip, width=\textwidth]{../R/boatshape-defense-1}};} % width=6, height=4.5
\uncover<2>{\node {\includegraphics[trim = 70mm 52.5mm 70mm 65mm, clip, width=\textwidth]{../R/boatshape-defense-2}};} % width=6, height=4.5
\end{tikzpicture}
}%
%\end{center}%}
\note<2>[item]{Prior set with different posterior sets (for different data $n=8$, $n=16$)}
\end{frame}


\section{Common-cause Failure Modelling}


\frame{\frametitle{Common-Cause Failures}
\begin{tikzpicture}
\uncover<1->{
\node {\includegraphics[scale=0.4]{./graph/800px-Fukushima_I_by_Digital_Globe.jpg}};
\node at (0,-3.5) {\tiny Source: Wikimedia Commons, \url{http://commons.wikimedia.org/wiki/File:Fukushima_I_by_Digital_Globe.jpg}};
}
\uncover<2->{
%\draw[white,fill=white] (-5,1) rectangle (5,3);
\node at (0,2.6) {\parbox{\textwidth}{
 \begin{alertblock}{common-cause failure} 
 \emph{simultaneous failure of several redundant components\\ due to a common or shared root cause}
 (H{\o}yland \& Rausand, 1994)
 \end{alertblock}
}};
}
\uncover<3>{
\draw[white,fill=white] (-0.2,-3.3) rectangle (6,0.75);
\node at (3,-1) {\includegraphics[scale=0.3]{./graph/plant-systems-diesel-generator.jpg}};
\node at (3.05,-3.05) {\parbox{6.2cm}{\tiny Source: \url{http://www.diakont.com/solutions/nuclear-energy/plant-systems/diesel-generator-control-systems/}}};
}
\end{tikzpicture}

% Notes for the slide
\note<3>[item]{This is a recent, drastic example for common-cause failure.\\
So what is common-cause failure? \play\ }
\note<3>[item]{Here, at the Fukushima Dai-ichi Nuclear Power Plant,\\
this simultaneous failure refers to the emergency diesel generators. \play\ }
\note<3>[item]{They are there to power the emergency cooling systems when off-site power is lost.}
\note<3>[item]{As cooling after shutdown is so important, there are usually 2 -- 4 EDG per reactor,
each of which can provide sufficient energy for the cooling systems alone.}
\note<3>[item]{At Fukushima Dai-ichi, there were 12 generators for 6 reactors.}
\note<3>[item]{When the T\={o}hoku earthquake, the reactors were shut down automatically,
and the EDG started to produce power for cooling..}
\note<3>[item]{They stopped working simultaneously because the machine rooms in the basement
were flooded by the tsunami wave.}
\note<3>[item]{So, common-cause failure compromises the reliability of redundant systems,
with components not failing independently.}
\note<3>[item]{So this must be included in overall systems reliability analysis}
\note<3>[item]{This lesson was learned at another accident at a nuclear power station, \play\ }
}

\frame{
\frametitle{Common-Cause Failure Modelling}

\begin{columns}%[T]
\begin{column}{0.7\textwidth}\centering
\includegraphics[scale=1.2]{./graph/800px-Three_Mile_Island_nuclear_power_plant.jpg}
\end{column}
\begin{column}{0.35\textwidth}\centering
\includegraphics[scale=0.95]{./graph/Graphic_TMI-2_Core_End-State_Configuration.png}
\end{column}
\end{columns}
\vspace*{1ex}
\begin{tiny}
Above: CDC, \url{http://phil.cdc.gov/phil/} ID 1194\\[1ex]
Right: Wikimedia Commons,\\[-1.8ex]
\url{http://commons.wikimedia.org/wiki/File:Graphic_TMI-2_Core_End-State_Configuration.png}
\end{tiny}
% Notes for the slide
\note<1>[item]{Lesson learned at another accident:\\ Three Mile Island (1979) where also a meltdown happened (right)}
\note<1>[item]{So, realising that reliability through redundancy is jeopardised by common-cause failure,
models including common-cause failure were developed in the wake of the Three Mile Island accident.}
\note<1>[item]{The most popular model is called \emph{basic parameter model}.}
\note<1>[item]{We will look at a parameterisation of this model called\ldots}
}

\subsection{Alpha-Factor Model}

\begin{frame}{Alpha-Factor Model: Definition}
  \begin{block}{Alpha-Factor Model} %\cite{1988:mosleh::common:cause}]
    Multinomial distribution $\mult(\vec{n}\mid\vec{\alpha})$ for common-cause failures \\
    in a $k$-component system\vspace*{-2ex}
    %the number of redundant components in the common-cause component group
    \begin{align*}
      p(\vec{n}\mid\vec{\alpha})=\prod_{j=1}^k\alpha_j^{n_j}\\[-5ex]
    \end{align*}
    where
    \begin{itemize}
    \item \alert{alpha-factor}
      $\alpha_j\coloneqq$
      \parbox[t]{0.6\textwidth}{%
        probability of $j$ of the $k$ components \\
        failing due to a common cause \\
        given that failure occurs
      }
    \item \alert{failure count}
      $n_j\coloneqq$ corresponding number of failures observed
    \item $\vec{n}$ denotes $(n_1,\dots,n_k)$ and $\vec{\alpha}$ denotes $(\alpha_1,\dots,\alpha_k)$
    \end{itemize}
  \end{block}
  (the model actually serves to estimate failure \emph{rates},\\ but the above is all what matters in this talk)
%\note<1>[item]{Slide on inferences in alpha-factor model commented out}
\end{frame}

\iffalse
\begin{frame}{Alpha-Factor Model: Inference}
  \begin{block}{Inference}
    involves rational functions \\
    of probabilities $\vec{\alpha}$ of common-cause failures\\
    and of total failure rate $q_t$ for individual components
  \end{block}
  \begin{center}
    \textbf{this talk: focus on $\vec{\alpha}$ only}
  \end{center}
\end{frame}
\fi

\begin{frame}{Alpha-Factor Model: Parameter Estimation}

%\begin{itemize}
maximum likelihood estimator: 
    \begin{align*}
      \hat{\alpha}_j &= \frac{n_j}{n},\quad \text{where $\textstyle\sum_{j=1}^n n_j = n$}
    \end{align*}
%\end{itemize}

\uncover<2->{%
  \begin{alertblock}{The Problem}%{The Bad News}
    \begin{itemize}[<+->]
    \item typically, for $j\ge 2$, the $n_j$ are very low \\
      with zero being quite common for larger $j$
    \item zero counts = flat likelihoods \quad\then\quad $\hat{\alpha}_j =$?
%      standard techniques such as MLE can struggle \\
%      to produce sensible inferences for this problem
%    \item need to rely on \alert{epistemic information}\\
%       \quad \then\ Bayesian inference
    \end{itemize}
  \end{alertblock}}
\uncover<3->{%
  \begin{center}
    \textbf{\then\ need to rely on \alert{epistemic information}: Bayesian inference}
  \end{center}}
\end{frame}


\subsection{Bayesian Inference}

\begin{frame}{Bayesian Inference: Dirichlet Prior}
\uncover<1->{%
%  $\vec{\alpha}$ considered as uncertain parameter on which we put\dots
  \begin{block}{Dirichlet-Multinomial Model}
  \vspace*{-2.0ex}
  \begin{columns}
  \begin{column}{0.65\textwidth}
    \begin{align*}
      p(\vec{\alpha}\mid \nzg,\byzr) \propto \prod_{j=1}^k\alpha_j^{\nzg\yzjr-1}
    \end{align*}
  \end{column}  
  \begin{column}{0.35\textwidth}
    \rule{0ex}{3ex}where $(\nzg, \byzr)$\\ are \emph{hyperparameters}
 \end{column}
 \end{columns}\vspace*{-2.7ex}
    \begin{align*}
      \nzg &> 0\\[-1.1ex]
      \byzr &\in \Delta=\Big\{
        (\yzor,\dots,\yzkr)\colon \yzor\ge0,\dots,\yzkr\ge 0,\,\sum_{j=1}^k \yzjr=1
      \Big\}
    \end{align*}%
    \vspace*{-4.8ex}%
    \begin{align*}
      p(\vec{\alpha}\mid\nzg,\byzr,\vec{n})
      &=
      p(\vec{\alpha}\mid\nng,\bynr)
      =
      \prod_{j=1}^k\alpha_j^{\nng\ynjr-1}\rule{9ex}{0ex}
    \end{align*}
  \end{block}}
\uncover<2->{%
  \vspace*{-0.5ex}
  \begin{block}{Interpretation}
    \begin{itemize}
    \item $\byzr$ = \alert{prior expectation of $\vec{\alpha}$}, i.e., a prior guess for $\frac{n_j}{n}$, $j=1,\ldots,n$
    \item $\nzg$ = determines \alert{spread} and \alert{learning speed} %(see next slide)
    \end{itemize}
  \end{block}}
\note<1>[item]{draw 1d and 2d simplex on black board}
\end{frame}

\iffalse
\begin{frame}{Bayesian Inference: Dirichlet Posterior}
  \begin{itemize}[<+->]
  \item posterior density for $\vec{\alpha}$ is again Dirichlet \ \ \then\ \ \alert{conjugacy}\\
  \alert{update parameters}: $\nzg \to\ \nng$, $\byzr \to \bynr$
    \begin{align*}
      p(\vec{\alpha}\mid\nzg,\byzr,\vec{n})
      &=
      p(\vec{\alpha}\mid\nng,\bynr)
      \propto
%      \prod_{j=1}^k\alpha_j^{\nzg\yzjr+n_j-1}
      \prod_{j=1}^k\alpha_j^{\nng\ynjr-1}
    \end{align*}
\item posterior expectation of $\alpha_j$:
    \begin{align*}
      %\label{eq:dirichlet:posterior:predictive}
      %\colorbox{red!10!white}{$\displaystyle
      \E[\alpha_j\mid\nzg,\byzr,\vec{n}]
      &=
      \E[\alpha_j\mid\nng,\bynr] 
      =
      \int_{\Delta}\alpha_j p(\vec{\alpha}\mid\nzg,\byzr,\vec{n})\dd\vec{\alpha} \\
      &= \ynjr
      =
      %\frac{n_j+st_j}{N+s}
      %=
      %\textcolor{green!50!black}{\frac{N}{N+s}}
      %\textcolor{blue}{\frac{n_j}{N}}
      %+
      %\textcolor{green!50!black}{\frac{s}{N+s}}
      %\textcolor{blue}{t_j}
      \colorbox{lmugreen!10!white}{$\displaystyle
      \frac{\nzg}{\nzg+n} \cdot \yzjr + \frac{n}{\nzg+n} \cdot \frac{n_j}{n}
      $}
    \end{align*}
%    where $n=\sum_{j=1}^k n_j$ is total number of observations
  \end{itemize}
  %\vspace{1em}
  \uncover<3->{%
%   \begin{alertblock}{}%{we will focus on $\E[\alpha_j\mid\nng,\bynr]$}
  \begin{center}
    \textbf{we will focus on $\E[\alpha_j\mid\nng,\bynr]$} \\
    (in a decision context, this expectation would typically end up \\
    in expressions for expected utility)
  \end{center}}
%   \end{alertblock}}
\end{frame}
\fi

\begin{frame}{Bayesian Inference: Example}
  
  \begin{exampleblock}{Example (Kelly \& Atwood, 2011)}
    Consider a system with four redundant components ($k=4$). \\
    %The probability of
    %$j$ out of $k$ common-cause failures, given that failure has happend,
    %was denoted by $\theta_j$.
    The analyst specifies
    the following prior expectation $\mu_{\text{spec},j}$
    for each $\alpha_j$:
    \begin{align*}
      %\label{eq:example:muspec}
      \mu_{\text{spec},1}&=0.950
      &
      \mu_{\text{spec},2}&=0.030
      &
      \mu_{\text{spec},3}&=0.015
      &
      \mu_{\text{spec},4}&=0.005
    \end{align*}
    We have 36 observations, in which 35 showed one component failing,
    and 1 showed two components failing:
    \begin{align*}
      n_1&=35
      &
      n_2&=1
      &
      n_3&=0
      &
      n_4&=0
    \end{align*}
  \end{exampleblock}
\vspace*{-2ex}%
\uncover<2->{%
    \begin{align*}
      %\label{eq:dirichlet:posterior:predictive}
      %\colorbox{red!10!white}{$\displaystyle
      \text{\bf we focus on }
      \E[\alpha_j\mid\nng,\bynr] 
      %=
      %\int_{\Delta}\alpha_j p(\vec{\alpha}\mid\nzg,\byzr,\vec{n})\dd\vec{\alpha} \\
      = \ynjr
      =
      \frac{\nzg}{\nzg+n} \cdot \yzjr + \frac{n}{\nzg+n} \cdot \frac{n_j}{n}
    \end{align*}
    (in a decision context, this expectation would typically end up \\
    in expressions for expected utility)}
\end{frame}

\begin{frame}[label=noninf-back]{Example: Non-Informative Priors}
  \alert{large variation in posterior under different \hyperlink{noninf-app}{non-informative priors}}
  \begin{itemize}%[<+->]
  \item<1-> with constrained maximum entropy prior\\ (Atwood, 1996; Kelly \& Atwood, 2011): %\cite{1996:atwood,2011:kelly:atwood}):
    \begin{align*}
      \E[\alpha_1\mid\nng,\bynr]&=0.967
      &
      \E[\alpha_2\mid\nng,\bynr]&=0.028
      \\
      \E[\alpha_3\mid\nng,\bynr]&=0.003
      &
      \E[\alpha_4\mid\nng,\bynr]&=0.001
    \end{align*}
  \item<2->
  with uniform prior ($\yzjr=0.25$ and $\nzg=4$):
  \begin{align*}
    \E[\alpha_1\mid\nng,\bynr]&=0.9
    &
    \E[\alpha_2\mid\nng,\bynr]&=0.05
    \\
    \E[\alpha_3\mid\nng,\bynr]&=0.025
    &
    \E[\alpha_4\mid\nng,\bynr]&=0.025
  \end{align*}
  \item<3->
  with Jeffrey's prior ($\yzjr=0.25$ and $\nzg=2$):
  \begin{align*}
    \E[\alpha_1\mid\nng,\bynr]&=0.9342
    &
    \E[\alpha_2\mid\nng,\bynr]&=0.0395
    \\
    \E[\alpha_3\mid\nng,\bynr]&=0.0132
    &
    \E[\alpha_4\mid\nng,\bynr]&=0.0132
  \end{align*}
  \end{itemize}%
% The degree of variation in the posterior under different priors
  % is evidently somewhat alarming.
  % In the next section, we aim to robustify the model
  % by using sets of priors from the start.
\note<1>[item]{emphasize variation for $\alpha_4$, the crucial part of the model\\
(all components fail): $0.1\%$ vs $2.5\%$ vs $1.3\%$}
\end{frame}

\subsection{Imprecise Dirichlet Model}

\begin{frame}{Imprecise Dirichlet Model: Definition}
\uncover<1->{%
Troffaes, Walter \& Kelly (2013): model vague prior info more cautiously
\begin{block}{Imprecise Dirichlet Model (IDM) for Common-Cause Failure}
use a \alert{set of hyperparameters}
    % \cite[p.~224, \S 5.4.3]{1991:walley} \cite[p.~32, \S 6]{1996:walley::idm},
    %\cite{1991:walley,1996:walley::idm}
    (Walley 1991, 1996):
    \begin{align*}
      %\label{eq:hyperparams:boxmodel}
      %\mathcal{H}
      \PZc
      =
      \left\{
        (\nzg,\byzr)
        \colon
        \nzg\in[\nzlg,\nzug],\,
        \byzr\in\Delta,\,
        \yzjr\in[\yzjlr{j},\yzjur{j}]
      \right\}
    \end{align*}
\end{block}\vspace*{-0.5ex}}\uncover<2->{%
\begin{block}{Interpretation}
\begin{itemize}
  \item we are doing a \alert{sensitivity analysis} (\'a la robust Bayes)\\ over $(\nzg,\byzr) \in \PZc$
  \item we take a \alert{set of priors} %$\text{conv}\big(\big\{p(\alpha\mid\nzg,\byzr) \colon (\nzg,\byzr) \in \PZc\big\}\big)$
  based on $\PZc$ 
  as model for prior information %(details later)
%  (resp.\ all convex mixtures of them)
\end{itemize}
\end{block}}\uncover<3>{%
Analyst has to specify (`elicit') \\
    bounds $[\nzlg,\nzug]$ and
    bounds $[\yzjlr{j},\yzjur{j}]$ for each $j\in\{1,\dots,k\}$ }%\\
\end{frame}

\begin{frame}{Imprecise Dirichlet Model: Elicitation}
  \begin{itemize}[<+->]
  \item   $[\yzjlr{j},\yzjur{j}]$: Cautious interpretation of prior specifications $\mu_{\text{spec},j}$:
  \begin{align*}
    [\yzjlr{1},\yzjur{1}]&=[0.950,1]
    &
    [\yzjlr{2},\yzjur{2}]&=[0,0.030]
    \\
    [\yzjlr{3},\yzjur{3}]&=[0,0.015]
    &
    [\yzjlr{4},\yzjur{4}]&=[0,0.005]
  \end{align*}
  \item  $[\nzlg,\nzug]$: Good (1965): \\ %\cite{1965:good}: \\
    \begin{center}
    reason about posterior expectations for hypothetical data
    \end{center}
  \end{itemize}
  \uncover<3->{%
  \begin{alertblock}{}
    $\nzug$ = number of one-component failures required \\
    to reduce the upper probabilities of multi-component failure by half
  \end{alertblock}
  \begin{alertblock}{}
    $\nzlg$ = number of multi-component failures required \\
    to reduce the lower probability of one-component failure by half
  \end{alertblock}}
\note<3>[item]{choice of bounds for $\nzg$ is not so straight-forward,
so we took advice from Good and reason about posterior expectations for hypothetical data,
also known as pre-posterior analysis}
\note<3>[item]{These rules are a bit difficult to grasp at first,\\ so I'll explain them with the example:}
\end{frame}

\begin{frame}{Imprecise Dirichlet Model: Elicitation}
\uncover<1->{%
\begin{alertblock}{}
    $\nzug$ = number of one-component failures required \\
    to reduce the upper probabilities of multi-component failure by half
  \end{alertblock}
  \begin{alertblock}{}
    $\nzlg$ = number of multi-component failures required \\
    to reduce the lower probability of one-component failure by half
  \end{alertblock}}
  \uncover<2->{%
  Reasonable values in example:
  \begin{itemize}
  \item 
    $\nzug=10$: %\\
    after observing $10$ one-component failures \\
    \then\ halve upper probabilities of multi-component failures
  \item
    $\nzlg=1$: %\\
    immediate multi-component failure
    \\
    \then\ keen to reduce lower probability for one-component failure
  \end{itemize}}
  \uncover<3->{%
  \alert{Difference between $\nzlg$ and $\nzug$} reflects a \alert{level of caution}:\\
  The rate at which we reduce upper probabilities \\
  is less than the rate at which we reduce lower probabilities \\}
  %\then\ reflects a \alert{level of caution}
\end{frame}

\begin{frame}{Imprecise Dirichlet Model: Inference}
%  \begin{center}
%    prior bounds + likelihood $\to$ posterior bounds
%  \end{center}
%\vspace*{-2ex}
With $[\nzlg, \nzug] = [1, 10]$, we get\ldots
\begin{block}{prior bounds + data $\to$ posterior bounds}
  \begin{center}
  \begin{tabular}{c|cc|c|cc}
    $j$ & $\yzjlr{j}$ & $\yzjur{j}$ & $n_j$ & $\El[\alpha_j\mid\PNc]$ &$\Eu[\alpha_j\mid\PNc]$ \\
    \hline
    1 & 0.950           & 1\phantom{0.00} & 35           & 0.967\phantom{0} & 0.978\phantom{00}\\
    2 & 0\phantom{0.00} & 0.030           & \phantom{0}1 & 0.0270           & 0.0283\phantom{0}\\
    3 & 0\phantom{0.00} & 0.015           & \phantom{0}0 & 0\phantom{0.000}  & 0.00326\\
    4 & 0\phantom{0.00} & 0.005           & \phantom{0}0 & 0\phantom{0.000}  & 0.00109
  \end{tabular}
  \end{center}
\end{block} 
  \begin{itemize}%[<+->]
  \item<2-> \alert{Bounds}, rather than precise values, are desirable \\
    due to inferences being strongly sensitive to the prior \\
    particularly when faced with zero counts
  \item<3-> Simple ways to elicit the parameters of the model \\
    by \alert{reasoning on hypothetical data} \\
%    rather than by maximum entropy arguments
%  \item<4-> Is it possible to generalise this method to other problems?
  \end{itemize}
\note<3>[item]{So what does this choice of parameter set mean for our number example?
We now have a posterior model that is based on a clear behavioural choice of prior,
expressing our prior ideas about $\vec{\alpha}$ which is still flexible enough to lead to reasonable posterior inferences}
\end{frame}


\frame{\frametitle{Conclusion}

\begin{itemize}%[<+->]
\item<1-> Conjugate priors are a convenient tool for Bayesian inference\\ but there are some pitfalls
 \begin{itemize}%[<+->]
 \item Hyperparameters $\nzg,\yzr$ are easy to interpret and elicit
 \item Averaging property makes calculations simple, but leads to
 inadequate model behaviour in case of \pdc
 \end{itemize}
\item<2-> Sets of conjugate priors maintain advantages \& mitigate issues %find a sweet spot in between
 \begin{itemize}%[<+->]
 \item Hyperparameter set shape is important %a crucial modeling choice
 \item Reasonable choice: \emph{rectangular} $\PZc = [\nzlg, \nzug] \times [\yzlr, \yzur]$\\
 (Walter \& Augustin 2009: \emph{generalised iLUCK-models}, \hyperlink{luck-app}{\texttt{luck}})
 \item Bounds for hyperparameters are easy to interpret and elicit
 \item Additional imprecison in case of \pdc\\
 leads to \alert{cautious inferences if, and only if, caution is needed}
 \item Shape for more precision in case of strong prior-data agreement is in development
 (joint work with Frank Coolen and Mi\c{k} Bickis)
 \end{itemize}
\end{itemize}
}

\frame{\frametitle{Works the Thesis is Based on}

\nocite{Walter2009a-short,Walter2010a-long,Walter2011a,Walter2012b,Troffaes2013a-short,itip-statinf,luck-package}

\printbibliography[heading=none,keyword=defensethesis]

%\bibliographystyle{plainnat}
%\bibliography{../bib/eigene,../bib/itip-refs,../bib/other-refs}
}

\frame{\frametitle{Further References}

\nocite{1994:hoyland,1965:good,1991:walley,1996:walley::idm,2011:kelly:atwood,1996:atwood,2006:evans,2005:quaeghebeurcooman,%
1994:berger,2012:benavolizaffalon-short}

\printbibliography[heading=none,keyword=defensefurther]

%\bibliographystyle{plainnat}
%\bibliography{../bib/eigene,../bib/itip-refs,../bib/other-refs}
}

%---- Appendix: Foundational arguments for IP
\begin{frame}[label=foundational-app]{Arguments for Imprecise Bayesian Inference
\hyperlink{foundational-back<3>}{\beamerreturnbutton{}}}
\begin{itemize}
\item The formal arguments for the Bayesian approach as a coherent way of inferential reasoning
hold only if one can make \emph{``arbitrarily fine discriminations in judgement about unknowns and utilities''}
(Berger et al.\ 1994, p.~303). (Your subjective prior must \emph{exactly} express your preferences,
anticipating \emph{anything} that may happen.)
\item Walley's (1991) framework for coherent inference using imprecise probabilities
instead allows for incomplete and imprecise prior specifications.
%\item It properly reflects the knowledge base ***
\note<1>[item]{The unattainable precision required by usual Bayesian inference\\ is thus unnecessary (even as an ideal).}
\begin{itemize}
\item more realistic description of uncertainties that are often hidden by spuriously precise models (through `arbitrary' modeling decisions)
\item `overprecision' is often compensated %in statistical practice
by `taking models not too seriously' (``all models are wrong, but some are useful'')
\item why not use sensible, reliable models in the first place?
\end{itemize}
\end{itemize}
\end{frame}

\begin{frame}[label=hierarchical-app]{Hierarchical Modelling
\hyperlink{foundational-back<3>}{\beamerreturnbutton{}}}
\begin{itemize}
\item Usually: space of parametric priors indexed by parameter $\phi$.
\item Very useful if $\phi$ has a clear interpretation (e.g., as a global mean),\\
such that a prior on $\phi$ can be meaningfully elicited.
\item Noninformative priors on $\phi$ can be problematic\\ (incoherence, improper posteriors).
\begin{exampleblock}{Example (e.g., Walley 1991, p.~232)}
$x_i \stackrel{iid}{\sim} \norm(\mu_i,1)$, $\mu_i \stackrel{iid}{\sim} \norm(\mu,\sigma^2)$, $i=1,\ldots,n$.\\
%For the noninformative hyperprior
For $p(\mu,\sigma^2) \propto \sigma^{-1}$, the posterior is improper.
\end{exampleblock}
\item High-variance priors on $\phi$ do not express ignorance about $\phi$,\\
but a strong prior belief that $|\phi|$ is large (Walley 1991, p.~233).
\end{itemize}
\end{frame}


%---- Appendix: Updating and Mixture Commute

\begin{frame}[label=commute-app]{Updating and Mixture Commute \hyperlink{commute-back<3>}{\beamerreturnbutton{}}}

Let
\begin{align*}
p_m(\vartheta\mid \nz_1\!,\yz_1\!, \nz_2\!,\yz_2\!, \kappa) &:=
   \kappa\,  p(\vartheta\mid\nz_1\!,\yz_1) +
(1-\kappa)\, p(\vartheta\mid\nz_2\!,\yz_2) \,,
\end{align*}
with marginals
\begin{align*}
f_1(\x) &= \int_\Theta\! f(\x\mid\vartheta) p(\vartheta\mid \nz_1\!,\yz_1) \dd\vartheta \,, \\
f_2(\x) &= \int_\Theta\! f(\x\mid\vartheta) p(\vartheta\mid \nz_2\!,\yz_2) \dd\vartheta \,, \\
f_m(\x) &= \int_\Theta\! f(\x\mid\vartheta) p_m(\vartheta\mid \nz_1\!,\yz_1\!, \nz_2\!,\yz_2\!, \kappa) \dd\vartheta \\
        &=    \kappa  \int_\Theta f(\x\mid\vartheta) p(\vartheta\mid \nz_1\!,\yz_1) \dd\vartheta 
         + (1-\kappa) \int_\Theta f(\x\mid\vartheta) p(\vartheta\mid \nz_2\!,\yz_2) \dd\vartheta \\
%        &=    \kappa  \int_\Theta f_1(\x) p(\vartheta\mid \nn_1,\yn_1) \dd\vartheta 
%         + (1-\kappa) \int_\Theta f_2(\x) p(\vartheta\mid \nn_2,\yn_2) \dd\vartheta \\
        &= \kappa f_1(\x) + (1-\kappa) f_2(\x) \,.
\end{align*}

\end{frame}

\frame{\frametitle{Updating and Mixture Commute \hyperlink{commute-back<3>}{\beamerreturnbutton{}}}

%Then
\begin{align*}
\lefteqn{
p_m(\vartheta\mid \nz_1,\yz_1, \nz_2,\yz_2, \kappa, \x) } \hspace*{8ex} & \\[2ex]
 &= \frac{f(\x\mid\vartheta)}{f_m(\x)} \big(
     \kappa\,  p(\vartheta\mid\nz_1,\yz_1) +
  (1-\kappa)\, p(\vartheta\mid\nz_2,\yz_2) \big) \\
 &=  \kappa\,  \frac{f(\x\mid\vartheta) p(\vartheta\mid\nz_1,\yz_1)}{f_m(\x)} +
  (1-\kappa)\, \frac{f(\x\mid\vartheta) p(\vartheta\mid\nz_2,\yz_2)}{f_m(\x)} \\
 &=  \kappa\,  \frac{f_1(\x) p(\vartheta\mid\nn_1,\yn_1)}{f_m(\x)} +
  (1-\kappa)\, \frac{f_2(\x) p(\vartheta\mid\nn_2,\yn_2)}{f_m(\x)} \\
 &= p_m(\vartheta\mid \nn_1,\yn_1, \nn_2,\yn_2, \kappa^*)\,,
\end{align*}
\begin{align*}
\text{where}\hspace{5ex}
\kappa^* &= \kappa\, \frac{f_1(\x)}{f_m(\x)}
          = \frac{\kappa\, f_1(\x)}{\kappa\, f_1(\x) + (1-\kappa)\, f_2(\x)}\,. 
\end{align*}

}

\frame{\frametitle{Updating and Mixture Commute \hyperlink{commute-back<3>}{\beamerreturnbutton{}}}

\begin{itemize}
\item updated mixture distribution is a mixture of the updated components
with mixture parameter $\kappa^*$ instead of $\kappa$.
\item convex hull of prior components\\ $=$ set of prior mixture distributions with $\kappa \in [0,1]$.
\item for any $\kappa \in [0,1]$, the corresponding $\kappa^* \in [0,1]$.
\item in fact, $\{\kappa^* \mid \kappa \in [0,1]\} = [0,1]$.
\item set of updated mixture distributions with $\kappa \in [0,1]$\\
$=$ convex hull of updated components ($\kappa^* \in [0,1]$).
\item arbitrary number of components by complete induction.
\end{itemize}
}

\begin{frame}[label=othermodels-app]{Other Models Based on Sets of Priors
\hyperlink{othermodels-back<3>}{\beamerreturnbutton{3}} \hyperlink{othermodels-back<10>}{\beamerreturnbutton{10}}}
\begin{itemize}
\item Neighbourhood models
\begin{itemize}
\item set of distributions `close to' a central distribution $\p_0$
\item common in robust Bayesian approaches 
\item example: $\varepsilon$-contamination class:
$\{ \p : \p = (1-\varepsilon) \p_0 + \varepsilon \q, \q \in \mathcal{Q} \}$
\item not necessarily closed under Bayesian updating
\end{itemize}
\item Density ratio class / interval of measures 
\begin{itemize}
\item set of distributions by bounds for the density function $p(\vartheta)$:
\begin{align*}
\mathcal{M}_{l,u} = \big\{ p(\vartheta) :
\exists c \in \posreals: l(\vartheta) \le c p(\vartheta) \le u(\vartheta)\big\}
\end{align*}
\item posterior set is bounded by updated $l(\vartheta)$ and $u(\vartheta)$
\item $u(\vartheta)/l(\vartheta)$ is constant under updating\\
\quad\then\ size of the set does not decrease with $n$\\
\quad\then\ too vague posterior inferences
\end{itemize}
\end{itemize}
\end{frame}


\begin{frame}{Other Models Based on Sets of Priors
\hyperlink{othermodels-back<3>}{\beamerreturnbutton{3}} \hyperlink{othermodels-back<10>}{\beamerreturnbutton{10}}}
\begin{itemize}
\item Discrete models
\begin{itemize}
\item discretize the parameter space as $\Theta = \{\theta_j\}_{j \in \{1,\ldots,m\}}$
\item set of distributions by bounds for $p(\vartheta_j)$ (or for expectations)
\item posterior bounds determined by linear programming algorithm
\item more flexibility at the cost of computational complexity
\item no clear measure for weight of prior information as compared to $n$
\end{itemize}
\end{itemize}
\begin{center}
\begin{tikzpicture}%[transform canvas={scale=0.33}]
\pgftransformscale{0.4}
% 
\coordinate (e100) at (90:6);
\coordinate (e010) at (210:6);
\coordinate (e001) at (-30:6);
\draw [thick] (e100) -- (e010) -- (e001) -- cycle;
% Achsen
\draw [thick,-stealth] (e001) -- +( 90:1.5) node [right] {$\vartheta_1$};
\draw [thick,-stealth] (e100) -- +(210:1.5) node [left]  {$\vartheta_2$};
\draw [thick,-stealth] (e010) -- +(-30:1.5) node [right] {$\vartheta_3$};
% graue Linien
%
\foreach \zehntel in {0.1,0.2,0.3,0.4,0.5,0.6,0.7,0.8,0.9,1}
{
 % Richtung p1
 \draw[gray,densely dotted] (barycentric cs:e100=\zehntel,e010=0,e001=1-\zehntel) -- %
                            (barycentric cs:e100=\zehntel,e010=1-\zehntel,e001=0);   %
 % Richtung p2
 \draw[gray,densely dotted] (barycentric cs:e100=1-\zehntel,e010=\zehntel,e001=0) -- %
                            (barycentric cs:e100=0,e010=\zehntel,e001=1-\zehntel);   %
 % Richtung p3
 \draw[gray,densely dotted] (barycentric cs:e100=0,e010=1-\zehntel,e001=\zehntel) -- %
                            (barycentric cs:e100=1-\zehntel,e010=0,e001=\zehntel);   %
}
%\draw[fill=gray!40] (barycentric cs:e100=0.75,e010=0.125,e001=0.125) circle (5pt);
% prior set
\draw[fill=gray!40] (barycentric cs:e100=0.4,e010=0.4,e001=0.2) -- %
                    (barycentric cs:e100=0.4,e010=0.1,e001=0.5) -- %
                    (barycentric cs:e100=0.7,e010=0.1,e001=0.2) -- cycle;
% data point
%\fill[black] (barycentric cs:e100=0.75,e010=0.125,e001=0.125) circle (5pt);
\begin{scope}[xshift=13cm]
\coordinate (e100) at (90:6);
\coordinate (e010) at (210:6);
\coordinate (e001) at (-30:6);
\draw [thick] (e100) -- (e010) -- (e001) -- cycle;
% Achsen
\draw [thick,-stealth] (e001) -- +( 90:1.5) node [right] {$\vartheta_1$};
\draw [thick,-stealth] (e100) -- +(210:1.5) node [left]  {$\vartheta_2$};
\draw [thick,-stealth] (e010) -- +(-30:1.5) node [right] {$\vartheta_3$};
% graue Linien
%
\foreach \zehntel in {0.1,0.2,0.3,0.4,0.5,0.6,0.7,0.8,0.9,1}
{
 % Richtung p1
 \draw[gray,densely dotted] (barycentric cs:e100=\zehntel,e010=0,e001=1-\zehntel) -- %
                            (barycentric cs:e100=\zehntel,e010=1-\zehntel,e001=0);   %
 % Richtung p2
 \draw[gray,densely dotted] (barycentric cs:e100=1-\zehntel,e010=\zehntel,e001=0) -- %
                            (barycentric cs:e100=0,e010=\zehntel,e001=1-\zehntel);   %
 % Richtung p3
 \draw[gray,densely dotted] (barycentric cs:e100=0,e010=1-\zehntel,e001=\zehntel) -- %
                            (barycentric cs:e100=1-\zehntel,e010=0,e001=\zehntel);   %
}
\draw[fill=gray!40] (barycentric cs:e100=0.5,e010=0.4,e001=0.1) -- %
                    (barycentric cs:e100=0.3,e010=0.2,e001=0.5) -- %
                    (barycentric cs:e100=0.7,e010=0.1,e001=0.2) -- cycle;
\end{scope}
\end{tikzpicture}
\end{center}
\end{frame}


\begin{frame}[label=luck-app]{\textbf{R} package \texttt{luck} \hyperlink{luck-back<4>}{\beamerreturnbutton{}}}
\begin{itemize}
\item \texttt{S4} implementation of the general canonical prior parameter structure
with rectangular sets $\PZc = [\nzlg, \nzug] \times [\yzlr, \yzur]$
\item lean subclasses for concrete sample distributions\\
(currently implemented: scaled normal, exponential)
\item currently on R-Forge, to be submitted to CRAN
%{\small (issue: \texttt{S4} classes and \texttt{S3}-style documentation seem to be incompatible)} 
\end{itemize}
{\small \url{install.packages("luck", repos="http://R-Forge.R-project.org")}}
\end{frame}

\begin{frame}{\textbf{R} package \texttt{luck} \hyperlink{luck-back<4>}{\beamerreturnbutton{}}}
\vspace*{-5ex}
\hspace*{10ex}
\begin{tikzpicture}[class/.style={draw, rectangle split, rectangle split parts=3, %
                                  every text node part/.style={text centered}, %
                                  font=\ttfamily, %
                                  fill=gray!30, text width=45mm, thick}, %
                    point/.style={coordinate},
                    transform canvas={scale=0.8}]%every node/.style=draw]
\node [class] (luck) at (0,0) {LuckModel \nodepart{second} \gruen{n0}: matrix \\
                                         \rot{y0}: matrix \\
                                         data: LuckModelData
                                         \nodepart{third}
                                         show() \\
                                         plot() \\
                                         unionHdi() \\
                                         \quad \vdots};
\node [coordinate, below=0.4cm of luck] (p0) {};
\node [coordinate, left=0.4cm of p0] (p1) {};
\node [class, text width = 47mm] (scalednormal) [below=0.4cm of p1] {ScaledNormalLuckModel\phantom{p} %
                                                                        \nodepart{second} \phantom{pS}
                                                                        \nodepart{third} singleHdi()\phantom{p}\\};
\draw [stealth'-, thick] (luck.south) -- (p0) -- (p1) -- (scalednormal.north); %node [midway,right,draw=none] {inherits};
\node[class, text width=45mm] (poisson) [right=0.5cm of scalednormal] {ExponentialLuckModel %
                                                                        \nodepart{second} \phantom{pS}
                                                                        \nodepart{third} singleHdi()\phantom{p}\\};
\node[class, text width=35mm, double copy shadow={shadow xshift=1mm,shadow yshift=-1mm}] %
 (more) [right=0.5cm of poisson] {\phantom{pS}\ldots\phantom{pS} \nodepart{second} \phantom{pS}
                                                                 \nodepart{third} singleHdi()\phantom{pS}\\};
\node [coordinate, above=0.4cm of poisson] (p2) {};
\draw [thick] (p0) -- (p2) -- (poisson.north);
\node [coordinate, above=0.4cm of more] (p3) {};
\draw [thick] (p2) -- (p3) -- (more.north);
%\draw[->] (0,0) -- (1,0);
\begin{scope}[transform canvas={scale=0.75}]
%\draw[->] (0,0.5) -- (1,0.5);
%\node [coordinate] (lmd) [right=0.5cm of luck] {}; 
\node [class] (data) at (6.8,1.5) {LuckModelData \nodepart{second} tauN: matrix\\ rawData: matrix
                                                             \nodepart{third} show()};
\node [coordinate, below=0.4cm of data] (dp0) {};
\node [coordinate, left=0.4cm of dp0] (dp1) {};
\node[class, text width=37mm] (scalednormaldata) [below=0.4cm of dp1] {ScaledNormalData\phantom{p} \nodepart{second} \phantom{h}
                                                                                                   \nodepart{third} show()};
\node[class, text width=37mm] (poissondata) [right=0.5cm of scalednormaldata] {ExponentialData \nodepart{second} \phantom{h}
                                                                                               \nodepart{third} show()};
\node[class, text width=30mm, double copy shadow={shadow xshift=1mm,shadow yshift=-1mm}] %
 (moredata) [right=0.5cm of poissondata] {\phantom{pD}\ldots\phantom{pD} \nodepart{second} \phantom{h}
                                                                         \nodepart{third}  show()};
\draw [stealth'-, thick] (data.south) -- (dp0) -- (dp1) -- (scalednormaldata.north); %node [midway,right,draw=none] {inherits};
\node [coordinate, above=0.4cm of poissondata] (dp2) {};
\draw [thick] (dp0) -- (dp2) -- (poissondata.north);
\node [coordinate, above=0.4cm of moredata] (dp3) {};
\draw [thick] (dp2) -- (dp3) -- (moredata.north);
\end{scope}
\node [coordinate] (snake) at (2.75,1.025) {};
\draw [dashed, very thick, -stealth'] (2.5,0.35) .. controls +(right:4mm) and +(up:4mm)
                                      .. (snake) .. controls +(up:4mm) and +(left:8mm) .. (3.2,1.7);
\end{tikzpicture}
\end{frame}

\begin{frame}[label=boat-app]{Strong Prior-Data Agreement Modelling \hyperlink{boat-back<2>}{\beamerreturnbutton{}}}
\includegraphics[trim = 15mm 45mm 25mm 60mm, clip, width=\textwidth]{../R/boatshape-prior}
\end{frame}

\begin{frame}[label=noninf-app]{`Noninformative' Priors \hyperlink{noninf-back<3>}{\beamerreturnbutton{}}}
\begin{itemize}
\item noninformative priors can be useful as a technical device\\
when they hardly influence the posterior
\item they are often improper: problematic for testing and model selection
\item %in contrast to the name they do contain subtantial information:\\
they give a (precise) probability for any subset $A$ of the parameter space,
which seems incompatible with the notion of ignorance
\item they express \emph{indifference} instead of \emph{ignorance}
\item a set of priors with $P(A) = [0,1]$ is noninformative
\item model framework allows for near-noninformative sets of priors\\
%(Walley 1996; Benavoli \& Zaffalon 2012)
 \begin{itemize}
 \item IDM (Walley 1996): range of $\yzjr = (0,1)\ \forall j$
 \item Benavoli \& Zaffalon (2012):\\ range of $\yzr = (-\infty,+\infty)$,
 while $\nzug$ decreases with $\yzr$\\ %$\nzg$ bounded above
 (to avoid $\nzg |\yzr| = \infty$, i.e.\ vacuous posterior inferences)
 \end{itemize}
\end{itemize}
\end{frame}

%---- reset total frame number
\addtocounter{framenumber}{-11} 


\end{document}
